\documentclass[intern]{cgMA}
\usepackage{url}
\usepackage{hyperref}
\usepackage{float}
\usepackage{pdfpages}
\usepackage{amsmath}
\usepackage{amssymb}
\usepackage[english,ngerman]{babel}
\title{Konzeption und Entwicklung einer kamerabasierten Analyse von Taekwondo-Techniken in einer Trainings- und Lern-App}
\author{William Hong}
\zweitgutachter{Dr. rer. nat. habil. Sabine Bauer}
\zweitgutachterInfo{(Head of General University Sports)}
\externLogo{7.46cm}{logos/ExternLogoPlaceholder}
\externName{DIN: NewTechnologies}
\begin{document}
\maketitle
\clearpage
\section*{Erklärung zur Verwendung von KI-Werkzeugen}
Im Rahmen der Erstellung dieser Arbeit wurde das generative KI-Werkzeug
Claude unterstützend eingesetzt. Der Einsatz beschränkte sich auf die allgemeine
Recherche, den Diskurs konzeptioneller Ideen für den Aufbau des Systems,
Hilfestellungen zur programmatischen Herangehensweise an mathematische Probleme
und zur beispielhaften Anwendung des OpenCV-Frameworks sowie auf die
Generierung erster Programmcode-Entwürfe. Sämtliche generierte Inhalte wurden
eigenständig überprüft, überarbeitet und in den Gesamtkontext der Arbeit integriert.
Der Aufbau des Systems und dessen programmatische Umsetzung, Validierung
und Evaluation wurden eigenständig durchgeführt. Die Verantwortung für
diese wissenschaftliche Arbeit liegt ausschließlich beim Verfasser.
\clearpage
\begin{abstract}
Eine saubere technische Ausführung ist beim Erlernen von Taekwondo Kicks besonders wichtig. Gerade außerhalb des Trainings fehlt jedoch häufig die direkte Rückmeldung durch eine Trainerin oder einen Trainer. Diese Arbeit beschäftigt sich deshalb mit der Frage, wie sich die Ausführung ausgewählter Taekwondo Kicks anhand von Smartphone Aufnahmen automatisiert analysieren und regelbasiert bewerten lässt.

Dafür wurde eine webbasierte Anwendung entwickelt, die drei Taekwondo Techniken untersucht. Dazu gehören der \textit{Front Kick}, der \textit{Roundhouse Kick} und der \textit{Spinning Back Kick}. Die Bewegungen werden mithilfe der markerlosen \textit{Pose Estimation} von MediaPipe erfasst. Auf Grundlage der dabei erkannten \textit{Landmarks} werden Gelenkwinkel, Positionen und die Ausrichtung des Kickbeins relativ zum Becken bestimmt. Eine Zustandsmaschine ordnet die Bewegung verschiedenen Phasen zu. Die Bewertung erfolgt anschließend regelbasiert anhand vorher festgelegter Kriterien und Schwellenwerte. Da die ausgeführte Technik vor der Analyse ausgewählt wird, ist kein zusätzlich trainiertes Modell zur Technikerkennung erforderlich.

Der entwickelte Prototyp wurde mit 19 Taekwondo Trainierenden getestet. Dabei wurde eine hohe wahrgenommene Gebrauchstauglichkeit erreicht. Auf der \textit{System Usability Scale} wurde ein Mittelwert von 88,0 Punkten erzielt. Auch die Rückmeldungen zur Verständlichkeit und zum wahrgenommenen Nutzen der Anwendung fielen überwiegend positiv aus. Besonders die Hinweise zur Verbesserung der Technik wurden als hilfreich wahrgenommen. Gleichzeitig zeigte sich, dass die Anwendung zwar verständlich darstellt, welches Kriterium erfüllt oder nicht erfüllt wurde, der genaue Weg zur jeweiligen Punktzahl für einen Teil der Teilnehmenden jedoch nicht ausreichend nachvollziehbar war.

Die Ergebnisse zeigen, dass sich ausgewählte Bewegungsmerkmale der drei betrachteten Taekwondo Kicks mit Smartphone Videos und markerloser Pose Estimation erfassen und in einem regelbasierten Prototyp bewerten lassen. Gleichzeitig werden Grenzen bei Bewegungen mit starker Körperrotation sichtbar. Dies betrifft insbesondere den \textit{Spinning Back Kick}, bei dem Verdeckungen und die gewählte Rotationsmessung zu größeren Unsicherheiten führen können. Eine unabhängige Validierung der Messwerte und Bewertungen gegenüber einer Referenzmessung oder fachlichen Trainerurteilen wurde im Rahmen dieser Arbeit nicht durchgeführt. Die vorliegenden Ergebnisse erlauben daher keine Aussage darüber, wie genau die berechneten Punktwerte die tatsächliche technische Qualität eines Kicks abbilden.

\bigskip
\noindent\textbf{Schlagwörter:} Taekwondo, markerlose Pose Estimation, MediaPipe, Bewegungsanalyse, regelbasierte Bewertung, Webanwendung
\end{abstract}

\begin{otherlanguage}{english}
\begin{abstract}
Proper technique is an important part of learning Taekwondo kicks. Outside regular training, however, direct feedback from a coach is often not available. This thesis therefore investigates how selected Taekwondo kicks can be analyzed from smartphone recordings and evaluated using a rule based approach.

For this purpose, a web based application was developed to analyze three Taekwondo techniques. These are the \textit{Front Kick}, the \textit{Roundhouse Kick}, and the \textit{Spinning Back Kick}. The movements are captured using markerless pose estimation with MediaPipe. Based on the detected body landmarks, joint angles, positions, and the orientation of the kicking leg relative to the pelvis are calculated. A state machine is then used to divide the movement into different phases. The evaluation is performed using predefined criteria and threshold values. Since the technique is selected before the analysis, no additional trained model is required for automatic technique recognition.

The developed prototype was evaluated with 19 Taekwondo practitioners. The results showed a high level of perceived usability. The application achieved a mean score of 88.0 on the \textit{System Usability Scale}. Participants also rated the comprehensibility and perceived usefulness of the application positively. In particular, the suggestions for improving technique were considered helpful. At the same time, the results showed that while the application clearly indicates which criteria were met or not met, the exact way in which an individual score was produced was not sufficiently clear to some participants.

The results show that selected movement features of the three examined Taekwondo kicks can be extracted from smartphone videos using markerless pose estimation and evaluated within a rule based prototype. At the same time, limitations become apparent for movements involving substantial body rotation. This is particularly relevant for the \textit{Spinning Back Kick}, where occlusion and the chosen method for representing rotation can introduce greater uncertainty. The measurements and evaluations were not independently validated against a reference measurement or expert coach ratings as part of this study. Therefore, the present results do not allow a conclusion about how accurately the calculated scores represent the actual technical quality of a kick.

\bigskip
\noindent\textbf{Keywords:} taekwondo, markerless pose estimation, MediaPipe, motion analysis, rule-based assessment, web application
\end{abstract}
\end{otherlanguage}

\tableofcontents
\clearpage
\section{Einleitung}
\subsection{Motivation}
Im Bereich der Sportbiomechanik stellt die Analyse sportlicher Bewegungen ein wichtiges Werkzeug dar, um technikbestimmende Bewegungsmerkmale zu untersuchen und Rehabilitationsprozesse zu unterstützen \cite{Colyer2018}. Die hierfür etablierten Laborverfahren ermöglichen eine hohe Messgenauigkeit, sind aufgrund ihrer Anschaffungskosten und technischen Komplexität jedoch nur eingeschränkt für einen breiten Einsatz geeignet \cite{Edriss2025}. Hinzu kommen methodische Voraussetzungen, die einer routinemäßigen Anwendung für den alltäglichen Trainingsbetrieb im Wege stehen. So erfordert die Erfassung kinematischer Daten bei markerbasierten Verfahren die Anbringung von Markern am Körper sowie kontrollierte Aufnahmebedingungen und eine zeitintensive Nachbearbeitung \cite{Colyer2018}. Daraus ergibt sich ein zunehmendes Interesse an automatisierten, markerlosen Verfahren zur Bewegungsanalyse \cite{Colyer2018}.
Fortschritte im Bereich der künstlichen Intelligenz haben in den vergangenen Jahren zu einer rasanten Entwicklung von Verfahren zur 2D- und 3D-Modellierung des menschlichen Körpers geführt \cite{Edriss2025}. Frei verfügbare Frameworks wie OpenPose, MediaPipe und AlphaPose ermöglichen dabei die automatische Detektion von \textit{Keypoints} des Körpers in Videoaufnahmen und können mit handelsüblichen Kameras eingesetzt werden \cite{Edriss2025}. Auf Grundlage dieser \textit{Keypoints} lassen sich unter anderem kinematische Kenngrößen wie Gelenkwinkel aus Videoaufnahmen bestimmen. Dies wurde bereits in verschiedenen Studien zur Analyse sportlicher Bewegungen untersucht \cite{Edriss2025}. Auch die Verwendung eines handelsüblichen Smartphones als Aufnahmegerät wurde in diesem Zusammenhang untersucht \cite{Babouras2024}.
Ein Anwendungsfeld für eine solche automatisierte Bewegungsanalyse sind Kampfsportarten. Echeverria und Santos \cite{Echeverria2021} begründen das damit, dass Kampfsporttechniken auf einem System von Bewegungen beruhen, die vorgegeben oder zumindest klar begrenzt sind und den Gesetzen der Physik folgen \cite{Echeverria2021}. Die Ausführung solcher Techniken ist somit durch definierte Bewegungsmerkmale charakterisiert, anhand derer eine technisch saubere Ausführung beschrieben werden kann. Hinzu kommt, dass ein Treffer im Wettkampf unter anderem von der Ausführung der Technik abhängt und sich eine stabile Ausführung erst über wiederholtes Üben entwickelt \cite{Echeverria2021}.
Im Taekwondo stehen insbesondere verschiedene Fußtechniken im Mittelpunkt, deren Ausführung durch komplexe biomechanische Bewegungsabläufe gekennzeichnet ist \cite{YilmazAtes2025}. Punkte werden erzielt, wenn ein Kick präzise und mit ausreichender Kraft den Rumpf oder den Kopf des Gegners trifft \cite{Falco2009}. Dabei ist das Verständnis der biomechanischen Eigenschaften dieser Techniken sowohl für die Optimierung der sportlichen Leistung als auch für die Reduktion des Verletzungsrisikos relevant \cite{YilmazAtes2025}. Beispielsweise untersuchten Falco et al. \cite{Falco2009} den Einfluss der Ausgangsdistanz auf die Ausführung des \textit{Roundhouse Kicks}. Dabei zeigte sich ein Einfluss der Ausgangsdistanz auf die Ausführungszeit, während sich die Aufprallkraft bei erfahrenen Wettkämpfern nicht signifikant unterschied \cite{Falco2009}. Solche Untersuchungen verdeutlichen die Bedeutung der Erfassung und Bewertung biomechanischer Merkmale bei der Analyse von Taekwondo-Techniken.
Taekwondo-Techniken weisen darüber hinaus strukturelle Eigenschaften auf, die eine automatisierte Technikanalyse begünstigen. Insbesondere stellen einzelne Fußtechniken zeitlich begrenzte Bewegungsabläufe dar, die sich anhand charakteristischer Bewegungsphasen beschreiben lassen. So wurde beispielsweise der Taekwondo \textit{Roundhouse Kick} in biomechanischen Untersuchungen in mehrere aufeinanderfolgende Phasen unterteilt, die durch charakteristische Bewegungsereignisse definiert sind \cite{Gavagan2017}. Die zeitliche und strukturelle Gliederung einer Technik ermöglicht es, spezifische technische Merkmale einzelnen Bewegungsphasen zuzuordnen und auf dieser Grundlage zu bewerten. Für eine regelbasierte Technikanalyse bietet diese Struktur somit eine geeignete Grundlage, da für einzelne Bewegungsphasen spezifische Bewertungskriterien formuliert werden können.
\subsection{Problemstellung}
Im Trainingsalltag stößt die Technikkorrektur allerdings an praktische Grenzen. Eine einzelne Trainerin oder ein einzelner Trainer kann im Gruppentraining nicht alle Schülerinnen und Schüler gleichzeitig im Blick haben, sodass ein großer Teil der Wiederholungen ohne Rückmeldung bleibt. Wang et al. \cite{Wang2026} beschreiben für das Taekwondo-Training, dass eine Lehrkraft in begrenzter Zeit 30 bis 40 Schülerinnen und Schüler betreut und unter diesen Bedingungen nur an wenigen Stellen verbale Rückmeldung geben kann \cite{Wang2026}. Außerhalb betreuter Trainingseinheiten steht den Trainierenden hingegen häufig kein unmittelbares fachliches Feedback zur Verfügung. Mündliches Feedback bezieht sich zudem auf die jeweilige Ausführung und wird ohne zusätzliche Dokumentation nicht dauerhaft festgehalten. Dadurch ist eine spätere Nachverfolgung der technischen Entwicklung erschwert.
Ein möglicher Ansatz zur Selbstkontrolle ist die Nutzung eines Spiegels in den Trainingshallen. Dieser ermöglicht zwar unmittelbar während der Ausführung eine visuelle Rückmeldung, setzt jedoch voraus, dass die Ausführenden diese Rückmeldung während der Bewegung wahrnehmen und interpretieren können. Eine Videoaufnahme kann hier weiterhelfen, da sich der Bewegungsablauf im Nachhinein in Ruhe und wiederholt betrachten lässt. Dass videobasiertes visuelles Feedback das motorische Lernen unterstützen kann, ist für den Sportunterricht belegt \cite{Moedinger2022}. Eine Videoaufnahme dokumentiert zwar die tatsächliche Bewegungsausführung, liefert jedoch nicht automatisch eine Bewertung ihrer technischen Qualität. Für eine eigenständige Beurteilung ist daher zunächst Wissen darüber erforderlich, anhand welcher Merkmale die technische Ausführung bewertet werden kann.
Ein möglicher Weg besteht darin, die Bewertung der Ausführung zu automatisieren. Systeme dieser Art existieren für einzelne Bewegungsfelder bereits \cite{Pham2022}. El-Rajab et al. \cite{ElRajab2025} weisen jedoch darauf hin, dass die automatische Bewertung der Bewegungsqualität bislang nur begrenzt untersucht wurde \cite{ElRajab2025}. Für Taekwondo liegt seit Kurzem ein erster Ansatz vor, der drei Grundtechniken aus Smartphone-Videos bewertet, die Erkennung der Technik jedoch einem neuronalen Netz überlässt \cite{Wang2026}. Ob und in welchem Umfang sich Drehtechniken mit ausgeprägter Körperrotation auf Grundlage von Smartphone-Videos automatisiert hinsichtlich ihrer technischen Ausführung bewerten lassen, ist bislang jedoch weitgehend offen.
\subsection{Zielsetzung}
Ziel dieser Arbeit ist die Konzeption und prototypische Umsetzung eines Systems, das einen mit einem Smartphone aufgenommenen Taekwondo Kick mithilfe markerloser \textit{Pose Estimation} analysiert und anhand festgelegter Kriterien regelbasiert bewertet. Betrachtet werden drei Techniken mit unterschiedlichem Rotationsanteil: der \textit{Front Kick}, der \textit{Roundhouse Kick} und der \textit{Spinning Back Kick}. Die zu analysierende Technik wird vor der Auswertung vom Nutzer ausgewählt. Dadurch ist kein zusätzlich trainiertes Modell zur automatischen Technikerkennung erforderlich. Die Bewertung basiert stattdessen auf aus den Körperlandmarken abgeleiteten Messgrößen und festgelegten Schwellenwerten.

Die Arbeit verfolgt dabei zwei miteinander verbundene Forschungsfragen:

\begin{enumerate}
    \item Wie lassen sich ausgewählte Merkmale der drei betrachteten Taekwondo-Kicks aus Smartphone-Aufnahmen mithilfe markerloser \textit{Pose Estimation} ableiten und regelbasiert bewerten, und welche technischen sowie konzeptionellen Grenzen zeigen sich dabei insbesondere bei ausgeprägter Körperrotation?
    \item Wie beurteilen Taekwondo Trainierende die Gebrauchstauglichkeit, Verständlichkeit, Vertrauenswürdigkeit und den wahrgenommenen Nutzen des entwickelten Prototyps?
\end{enumerate}

Die verwendeten Bewertungsgrenzen wurden aus vorhandener Technikliteratur, der Trainingspraxis und der prototypischen Entwicklung abgeleitet. Sie wurden im Rahmen dieser Arbeit nicht anhand einer unabhängigen Referenzmessung oder anhand systematischer Trainerurteile empirisch kalibriert. Für jedes Bewertungskriterium werden zwei Grenzwerte definiert, die den Bereich zwischen 0 und 100 Prozent festlegen. Ein Kriterium gilt in der Benutzeroberfläche ab einem berechneten Kriterienwert von 50 Prozent als bestanden.

Die Aufnahme erfolgt mit einem Smartphone aus einer seitlichen oder um etwa 45 Grad versetzten Perspektive. Pro Video wird eine einzelne Kickausführung analysiert. Neben der Kickanalyse umfasst der Prototyp eine ergänzende Vokabelfunktion zur Unterstützung der Vorbereitung auf Taekwondo-Gürtelprüfungen. Diese Funktion ist nicht Gegenstand der beiden Forschungsfragen und wird daher nur ergänzend betrachtet.

\subsection{Überblick der Arbeit}
Die Arbeit gliedert sich in sechs Kapitel. Kapitel 2 beschreibt die fachlichen und technischen Grundlagen. Dazu gehören die drei betrachteten Kicktechniken, markerlose Bewegungsanalyse, \textit{Pose Estimation} und grundlegende Ansätze zur automatisierten Bewegungsbewertung. Kapitel 3 ordnet die Arbeit gegenüber verwandten Ansätzen ein. Kapitel 4 beschreibt die Konzeption und Implementierung des Prototyps. Dabei werden die Verarbeitung der Körperlandmarken, die Phasenerkennung, die Berechnung der Messgrößen, die Bewertungskriterien und die Benutzeroberfläche erläutert. Kapitel 5 untersucht im Rahmen einer Nutzerstudie, wie die Zielgruppe die Gebrauchstauglichkeit, Verständlichkeit, das Vertrauen und den wahrgenommenen Nutzen der Anwendung bewertet. Kapitel 6 führt die Ergebnisse im Hinblick auf die Forschungsfragen zusammen, diskutiert die Grenzen des Ansatzes und leitet mögliche Weiterentwicklungen ab.

\section{Grundlagen}
\subsection{Taekwondo}
Taekwondo ist eine koreanische Kampfsportart, die sich seit der zweiten Hälfte des 20. Jahrhunderts zu einer international verbreiteten Sportart entwickelte \cite{Fong2011}. Seit den Olympischen Sommerspielen 2000 in Sydney ist Taekwondo Bestandteil des olympischen Wettkampfprogramms. Abbildung \ref{fig:taekwondo} zeigt eine Wettkampfsituation bei den Olympischen Spielen. Im internationalen Wettkampfsport werden insbesondere die Disziplinen Kyorugi und Poomsae unterschieden. Im Vollkontaktwettkampf (Kyorugi) treten zwei Athletinnen oder Athleten gegeneinander an, wobei Punkte durch Fuß- und Fausttechniken auf definierte Trefferflächen erzielt werden \cite{WTKyorugi2026}. Im Formenlauf (Poomsae) wird dagegen eine festgelegte Abfolge von Techniken ohne direkten Gegner ausgeführt und von Kampfrichtern bewertet \cite{WTPoomsae2024}.
\begin{figure}[htbp]
    \centering
    \includegraphics[width=0.8\textwidth]{Images/taekwondoolympics.jpg}
    \caption{Taekwondo-Wettkampf Olympische Sommerspiele 2016 \cite{WikiBild2026}}
    \label{fig:taekwondo}
\end{figure}
Neben dem Wettkampfsport spielt die technische Ausführung auch im regulären Training und bei Gürtelprüfungen eine zentrale Rolle. Im Training werden Grundtechniken wiederholt und sowohl ohne Ziel als auch gegen Trainingsgeräte wie Pratzen ausgeführt. Die Korrektur erfolgt dabei in der Regel durch die Trainerin oder den Trainer auf Grundlage der beobachteten Bewegungsausführung. Eine quantitative Erfassung der zugrunde liegenden Bewegungsmerkmale findet dabei nicht zwangsläufig statt.
Der Formenlauf ist zudem Bestandteil der Gürtelprüfungen. Neben praktischen Prüfungsinhalten umfasst eine Gürtelprüfung auch einen theoretischen Teil \cite{KukkiwonRules2018}. Dazu gehört unter anderem das Erlernen koreanischer Fachbegriffe für Techniken und Kommandos. Die sportliche Entwicklung ist über ein Gurtsystem strukturiert. Bis zum Schwarzgurt durchlaufen die Trainierenden verschiedene Kup-Grade, auf die die Meistergrade Dan folgen \cite{KukkiwonRules2018}. Der jeweilige Grad wird im Rahmen einer Prüfung erworben, in der sowohl praktische als auch theoretische Kenntnisse überprüft werden.
Die biomechanische Untersuchung von Taekwondo-Techniken ist sowohl für die Optimierung der sportlichen Leistung als auch für die Reduktion des Verletzungsrisikos relevant \cite{YilmazAtes2025}. Die vorhandene Forschung verteilt sich jedoch nicht gleichmäßig auf die verschiedenen Kicktechniken. Von den 86 Studien, die in einem systematischen Review zur Biomechanik von Taekwondo-Kicks berücksichtigt wurden, untersuchten 49 den Roundhouse Kick, der in dem Review als \textit{Dollyo Chagi} bezeichnet wird, während der Front Kick, dort als \textit{Ap Chagi} bezeichnet, in neun Studien betrachtet wurde. Der Spinning Back Kick, im Review als \textit{Dwit Chagi} bezeichnet, wird nicht als eigenständige Technik ausgewiesen, sondern der Kategorie der übrigen Kicktechniken zugeordnet, auf die zwölf Studien entfielen \cite{YilmazAtes2025}.
Für die vorliegende Arbeit sind sowohl Kyorugi als auch Poomsae relevant, jedoch aus unterschiedlichen Gründen. Im Kyorugi besitzen Fußtechniken eine besondere Bedeutung für die Punktevergabe. So sind Kopftreffer ausschließlich mit Fußtechniken zulässig und Kicks werden höher bewertet als Fausttechniken \cite{WTKyorugi2026}. Im Poomsae steht hingegen die Ausführung der Techniken selbst im Mittelpunkt, da keine Treffer gegen einen Gegner erzielt werden \cite{WTPoomsae2024}. Beide Disziplinen verdeutlichen damit die Bedeutung der technischen Ausführung von Kicks im Taekwondo.
Für die vorliegende Arbeit werden drei Techniken betrachtet, die sich insbesondere hinsichtlich ihrer Bewegungsstruktur und des Ausmaßes der Körperrotation unterscheiden: der Front Kick, der Roundhouse Kick und der Spinning Back Kick. Die folgenden Abschnitte beschreiben deren technische und biomechanische Merkmale, die anschließend als Grundlage für die automatisierte Bewertung dienen.
\subsection{Die untersuchten Techniken}
Die drei untersuchten Techniken unterscheiden sich vor allem hinsichtlich ihrer Bewegungsrichtung und des Ausmaßes der Körperrotation. Dennoch besitzen sie eine vergleichbare Grundstruktur aus Vorbereitung, Anheben beziehungsweise Anwinkeln des Kickbeins, Streckung und Rückführung. Gavagan und Sayers \cite{Gavagan2017} beschreiben für den \textit{Roundhouse Kick} die Abschnitte \textit{Preparation}, \textit{Chamber}, \textit{Extension} und \textit{Recoil} sowie das anschließende Aufsetzen des Beins. Dieses Literaturmodell dient als fachliche Orientierung für die spätere Phasenerkennung.

Für die Implementierung wird der Bewegungsablauf bewusst vereinfacht und auf die vier Zustände \texttt{idle}, \texttt{chamber}, \texttt{kick} und \texttt{rechamber} abgebildet. \texttt{idle} beschreibt dabei den Ausgangszustand vor dem Kick und die Rückkehr in diesen Zustand nach Abschluss der Bewegung. \texttt{chamber} steht für das Anheben und Anwinkeln des Kickbeins, \texttt{kick} für die Streckungsphase und \texttt{rechamber} für den Rückzug. Beim \textit{Spinning Back Kick} wird zusätzlich ein Zustand \texttt{rotation} verwendet, um die vorgelagerte Eindrehung des Körpers abzubilden. Die Softwarezustände sind damit keine unveränderte Übernahme des Literaturmodells, sondern eine für die regelbasierte Erkennung gewählte Abstraktion.

\subsubsection{Front Kick}
Der \textit{Front Kick}, im Taekwondo als \textit{Ap Chagi} bezeichnet, ist ein geradlinig nach vorne ausgeführter Fußstoß und zählt zu den grundlegenden Techniken des Taekwondo. Aus der Kampfstellung wird zunächst das Knie des Kickbeins ungefähr auf Hüfthöhe angehoben, während der Unterschenkel angewinkelt bleibt. Aus dieser Position wird der Unterschenkel nach vorne gestreckt. Als Trefferfläche dient dabei der Fußballen. Anschließend wird das Bein wieder angewinkelt und zurückgezogen, bevor es abgesetzt wird \cite{Gulia2022}.

Gulia und Shaw unterscheiden daneben den \textit{Front Toe Kick}, bei dem die Zehen als Trefferfläche eingesetzt werden \cite{Gulia2022}. Die Trefferfläche selbst wird in der vorliegenden Arbeit nicht ausgewertet, da die Landmarken an Ferse und Zehen nicht in die Bewertung eingehen (Abschnitt~4.2.2). Die Anwendung unterscheidet deshalb nicht automatisiert zwischen Varianten, die sich vor allem über die Fußstellung beziehungsweise Trefferfläche voneinander abgrenzen.

Für die vorliegende Analyse sind insbesondere die Beugung des Kickbeins in der \textit{Chamber} Phase, die anschließende Streckung und die Rückführung relevant. Zusätzlich wird berücksichtigt, wie stark das Knie während der Kick-Phase absinkt und wie das Kickbein im Treffmoment relativ zum Becken ausgerichtet ist. Abbildung~\ref{fig:ap_chagi} zeigt die Abfolge des Front Kicks.

\begin{figure}[h]
    \centering
    \includegraphics[width=0.8\textwidth]{Images/ap\_chagi.jpg}
    \caption{Phasen des \textit{Front Kicks} \cite{Gulia2022}}
    \label{fig:ap_chagi}
\end{figure}

\subsubsection{Roundhouse Kick}
Für die in dieser Arbeit betrachtete Technik wird durchgängig die englische Bezeichnung \textit{Roundhouse Kick} verwendet. In der koreanischen Taekwondo Terminologie wird die Technik als \textit{Dollyo Chagi} bezeichnet. In Teilen der biomechanischen Forschung findet sich daneben die Bezeichnung \textit{Bandal Chagi} \cite{Falco2009}. Um Mehrdeutigkeiten zwischen unterschiedlichen Quellen zu vermeiden, wird im weiteren Verlauf ausschließlich die Bezeichnung \textit{Roundhouse Kick} verwendet.

Der \textit{Roundhouse Kick} ist eine seitlich beziehungsweise bogenförmig ausgeführte Fußtechnik. Aus der Ausgangsstellung wird das Kickbein zunächst angehoben und das Knie angewinkelt. Anschließend wird das Knie gestreckt und der Fuß in Richtung des Ziels geführt. Gavagan und Sayers \cite{Gavagan2017} beschreiben dabei eine Vorbereitung, eine \textit{Chamber} Phase, die Streckung zum Ziel und den anschließenden Rückzug. Dieses Phasenmodell wird in der vorliegenden Implementierung auf die Zustände \texttt{idle}, \texttt{chamber}, \texttt{kick} und \texttt{rechamber} abgebildet (Abschnitt~2.2).

Ein wesentliches Merkmal des \textit{Roundhouse Kicks} ist die Rotation des Beckens um die Körperlängsachse. Gavagan und Sayers beschreiben die schnelle Beckenrotation als wichtigen Bestandteil der Technik \cite{Gavagan2017}. Sie unterstützt die bogenförmige Bewegungsbahn des Kickbeins und verändert die Ausrichtung des Körpers gegenüber dem Ziel. Für die technische Beurteilung sind daher insbesondere das Anwinkeln und Anheben des Kickbeins, die anschließende Streckung, die Rückführung sowie die Ausrichtung von Kickbein und Becken im Treffmoment relevant.

Die Trefferfläche des Fußes wird nicht berücksichtigt, da die Fußlandmarken in der vorliegenden Umsetzung nicht ausgewertet werden. Abbildung~\ref{fig:bandal_chagi} zeigt die Abfolge der Bewegungsphasen des Roundhouse Kicks.

\begin{figure}[h]
    \centering
    \includegraphics[width=0.8\textwidth]{Images/bandal\_chagi.jpg}
    \caption{Phasen des Roundhouse Kicks \cite{TaekwonWikiaDol}}
    \label{fig:bandal_chagi}
\end{figure}

\subsubsection{Spinning Back Kick}
Der \textit{Spinning Back Kick}, im Taekwondo als \textit{Dwit Chagi} bezeichnet, ist eine Rückwärtstritttechnik mit ausgeprägter Rotation um die Körperlängsachse. Aus der Ausgangsstellung wird der Körper zunächst eingedreht. Dabei wird das hintere Bein vom Boden gelöst und das Körpergewicht auf das Standbein verlagert. Pędzich et al. beschreiben für den Dwit Chagi eine Drehung um etwa 180 Grad auf dem Fußballen des vorderen Standbeins. Anschließend wird das Kickbein angehoben und in Richtung des Ziels gestreckt. Als Trefferflächen werden Ferse und Fußsohle angegeben \cite{pedzich2006}.

Wąsik und Góra beschreiben ebenfalls eine Rotation um etwa 180 Grad mit anschließender Anhebung und Streckung des Kickbeins. Nach dem Treffer wird das Bein zurückgeführt und eine Endposition eingenommen \cite{wasik2016}.

Für die technische Beurteilung sind damit die Eindrehung des Körpers, das Anwinkeln des Kickbeins während der Anhebung, die anschließende Streckung und die Rückführung relevant. Die vorliegende Implementierung kann die globale Körperrotation jedoch nur eingeschränkt erfassen. Das im Treffmoment verwendete Winkelmaß beschreibt in erster Linie die Ausrichtung des Kickbeins relativ zum Becken. Ergänzend wird eine aufsummierte Änderung der Beckenrichtung während des Bewegungsablaufs geführt. Die Grenzen beider Größen werden in Kapitel~6 diskutiert.

Die Trefferfläche wird auch bei dieser Technik nicht bewertet, da die Landmarken an Ferse und Zehen nicht ausgewertet werden.

Aus Sicht der videobasierten Analyse stellt der \textit{Spinning Back Kick} eine besondere Herausforderung dar. Durch die starke Rotation können einzelne Körperbereiche zeitweise von anderen Körperteilen verdeckt werden. Dies kann die Schätzung einzelner Landmarken beeinträchtigen. Wang et al. berichten bei einer Taekwondo-Technik mit ausgeprägter Hüftrotation ebenfalls von Selbstverdeckungen an Hüfte, Knie und Sprunggelenk als möglicher Ursache für eine geringere Erkennungsleistung \cite{Wang2026}. Da der \textit{Spinning Back Kick} eine noch weitergehende Körperdrehung beinhaltet, ist eine erhöhte Anfälligkeit für solche Effekte zu erwarten.

Abbildung~\ref{fig:dwit_chagi} zeigt die Abfolge der Bewegungsphasen des Spinning Back Kicks.

\begin{figure}[h]
    \centering
    \includegraphics[width=0.8\textwidth]{Images/dwit\_chagi.jpg}
    \caption{Phasen des \textit{Spinning Back Kicks} \cite{TaekwonWikiaDwi}}
    \label{fig:dwit_chagi}
\end{figure}

\subsection{Markerlose Bewegungsanalyse}
\subsubsection{Videobasierte Verfahren}
Markerbasierte optoelektronische Systeme bestimmen die Position der Körperpunkte über am Körper angebrachte Marker und erreichen dabei eine hohe Messgenauigkeit. Die damit verbundenen Anforderungen an Ausstattung, Aufnahmebedingungen und Nachbearbeitung wurden bereits in Abschnitt~1.1 beschrieben \cite{Colyer2018}.
Markerlose Verfahren bestimmen die Position der Körperpunkte stattdessen unmittelbar aus dem Kamerabild. Damit entfällt nicht nur das Anbringen der Marker, sondern auch die Bindung an eine Laborumgebung. Die Aufnahme kann dadurch auch außerhalb kontrollierter Laborbedingungen erfolgen. Allerdings hängt die Qualität der erfassten Bewegungsdaten von den Aufnahmebedingungen ab. Insbesondere Beleuchtung und Kameraposition können die Genauigkeit der Bewegungserfassung beeinflussen \cite{Thomas2026}.
Für die vorliegende Arbeit ist dieser Zusammenhang relevant, da die Aufnahmen nicht unter Laborbedingungen, sondern mit einem Smartphone im Trainingsumfeld entstehen. Die daraus abgeleiteten Anforderungen an die Aufnahme werden in Abschnitt~4.14 beschrieben.
\subsubsection{Pose Estimation}
Unter \textit{Pose Estimation} wird die automatische Bestimmung der Position definierter Körperpunkte in einem Bild verstanden. Diese Körperpunkte werden als \textit{Keypoints} oder \textit{Landmarks} bezeichnet. Erfasst werden vor allem Gelenke wie Schulter, Hüfte, Knie und Sprunggelenk, daneben aber auch Punkte an Gesicht und Händen \cite{Edriss2025}. Im weiteren Verlauf der Arbeit wird durchgängig die Bezeichnung \textit{Landmark} verwendet. Aus den so gewonnenen Punkten lässt sich ein Skelettmodell zusammensetzen, aus dem sich kinematische Kenngrößen wie Gelenkwinkel ableiten lassen \cite{Edriss2025}.
Die einzelnen Verfahren unterscheiden sich unter anderem in der Anzahl der \textit{Landmarks} und in der Form der Ausgabe. OpenPose erfasst 25 Punkte, das in MediaPipe enthaltene BlazePose 33 \cite{Edriss2025}. Werden die \textit{Landmarks} ausschließlich in Bildkoordinaten geliefert, handelt es sich um eine zweidimensionale \textit{Pose Estimation}. Verfahren, die zusätzlich eine Tiefeninformation schätzen, geben die Punkte als dreidimensionale Koordinaten aus \cite{Edriss2025}. Zweidimensionale Verfahren sind günstiger und einfacher einzurichten. Da ihnen die Tiefeninformation fehlt, erfassen sie Bewegungen, die aus der Aufnahmeebene herauslaufen, jedoch nur eingeschränkt und hängen stärker von der gewählten Perspektive ab. Als Beispiele nennen Edriss et al. \cite{Edriss2025} die Beugung und die Rotation eines Gelenks \cite{Edriss2025}. Für die vorliegende Arbeit ist diese Unterscheidung von Bedeutung, weil sämtliche Gelenkwinkel aus den dreidimensionalen Weltkoordinaten berechnet werden (Abschnitt~4.7).
\subsubsection{MediaPipe Pose Landmarker}
MediaPipe ist ein von Google entwickeltes \textit{Open-Source-Framework}, das eine Sammlung modularer Lösungen für Anwendungen des maschinellen Lernens bereitstellt \cite{MediaPipeTasks}. Erstmals vorgestellt wurde MediaPipe von Lugaresi et al. im Jahr 2019 \cite{Lugaresi2019}. Für die vorliegende Arbeit ist innerhalb dieser Sammlung insbesondere der \textit{Pose Landmarker} relevant. 
Der \textit{Pose Landmarker} arbeitet zweistufig. Zunächst detektiert ein Detektor die Anwesenheit einer Person im Bild und legt eine \textit{Region of Interest} (ROI) fest. Anschließend lokalisiert ein \textit{
Tracker} die \textit{Landmarks} innerhalb dieses Bildausschnitts und aktualisiert sie von Bild zu Bild. Dieses Vorgehen wird als Detector-Tracker-Prinzip bezeichnet \cite{Bazarevsky2020}. Der Vorteil liegt in der Rechenzeit, da die Detektion einer Person im Vollbild deutlich aufwendiger ist als das Tracking in einem zugeschnittenen Bildausschnitt. Der Detektor wird daher nur erneut ausgeführt, wenn die Person den Bildbereich verlässt oder die Verfolgung verloren geht \cite{Bazarevsky2020}.
Die Schnittstelle bildet diese Unterscheidung in drei Laufmodi ab. Der Modus \texttt{IMAGE} verarbeitet jedes Bild unabhängig, während \texttt{VIDEO} und \texttt{LIVE\_STREAM} die beschriebene Verfolgung nutzen und deshalb für jedes Bild einen Zeitstempel erwarten \cite{mediapipepose}.
Das Modell liefert 33 \textit{Landmarks}, die in Abbildung \ref{fig:pose_landmarks} und Tabelle \ref{tab:pose_landmarks} aufgeführt sind. Für jede \textit{Landmark} werden sowohl Bildkoordinaten als auch dreidimensionale Weltkoordinaten in Metern bereitgestellt. Zusätzlich gibt das Modell für jede \textit{Landmark} einen Sichtbarkeitswert aus, der angibt, wie wahrscheinlich der Punkt im Bild sichtbar ist \cite{mediapipepose}. Zugrunde liegt ein \textit{Convolutional Neural Network}, dessen Aufbau sich am Architekturprinzip von MobileNetV2 orientiert \cite{mobilenetv2, mediapipepose}.
Google bietet den \textit{Pose Landmarker} in den drei Varianten \textit{Lite, Full und Heavy} an, die sich in der Größe des verwendeten Modells und damit im Verhältnis von Genauigkeit und Laufzeit unterscheiden \cite{mediapipepose}. Welche Variante in der vorliegenden Arbeit eingesetzt wird, beschreibt Abschnitt~4.2.1.
Die Weltkoordinaten beziehen sich nicht auf den Bildursprung, sondern auf den Mittelpunkt zwischen beiden Hüften, und werden in Metern angegeben \cite{mediapipepose}. Sie sind damit hinsichtlich der Translation bereits normiert, sodass eine Verlagerung der Person im Raum die Koordinaten nicht verändert. Umgekehrt folgt daraus, dass die Weltkoordinaten die Bewegung ausschließlich relativ zum Körper beschreiben. Diese Eigenschaft wird in Abschnitt~4.3.3 bei der Normierung genutzt und wirkt sich zugleich auf die Geschwindigkeitsmessung aus (Abschnitt~4.11).
\begin{figure}[h]
    \centering
    \includegraphics[width=0.8\textwidth]{Images/pose\_landmarks\_index.png}
    \caption{Modell für Landmarken für Posen mit 33 Körpermarkierungen}
    \label{fig:pose_landmarks}
\end{figure}
\begin{table}[h]
\centering
\small
\begin{tabular}{c l c l}
\hline
\textbf{ID} & \textbf{Landmark} & \textbf{ID} & \textbf{Landmark} \\ \hline
0  & Nose              & 17 & Left pinky \\
1  & Left eye (inner)  & 18 & Right pinky \\
2  & Left eye          & 19 & Left index \\
3  & Left eye (outer)  & 20 & Right index \\
4  & Right eye (inner) & 21 & Left thumb \\
5  & Right eye         & 22 & Right thumb \\
6  & Right eye (outer) & 23 & Left hip \\
7  & Left ear          & 24 & Right hip \\
8  & Right ear         & 25 & Left knee \\
9  & Mouth (left)      & 26 & Right knee \\
10 & Mouth (right)     & 27 & Left ankle \\
11 & Left shoulder     & 28 & Right ankle \\
12 & Right shoulder    & 29 & Left heel \\
13 & Left elbow        & 30 & Right heel \\
14 & Right elbow       & 31 & Left foot index \\
15 & Left wrist        & 32 & Right foot index \\
16 & Right wrist       &    &  \\ \hline
\end{tabular}
\caption{Definition der 33 MediaPipe Pose Landmarks}
\label{tab:pose_landmarks}
\end{table}
\subsubsection{Automatisierte Bewegungsbewertung}
Ein Verfahren zur \textit{Pose Estimation} liefert ausschließlich Positionen. Ob eine Ausführung technisch korrekt war, ergibt sich daraus nicht von selbst. Zwischen den erfassten \textit{Landmarks} im Bild und einer Aussage über die Qualität der Ausführung liegt ein eigener Verarbeitungsschritt, in dem festgelegt wird, welche Größen aus den Positionen berechnet werden, welche Werte als korrekt gelten und in welcher Form das Ergebnis zurückgemeldet wird.
Ein Review von Tharatipyakul et al. \cite{Tharatipyakul2024} zu videobasiertem Bewegungsfeedback unterscheidet drei Wege, eine Bewegung zu bewerten \cite{Tharatipyakul2024}. Der erste rechnet die Bewertung über eine mathematische Formel aus den Messwerten aus. Der zweite ist regelbasiert und prüft die Ausführung gegen festgelegte Bedingungen und Schwellenwerte. Der dritte überlässt die Bewertung einem trainierten Modell. Der eigentlichen Bewertung ist in einigen Arbeiten ein Vorverarbeitungsschritt vorgelagert, der die Aufnahme zeitlich auf eine Referenzbewegung ausrichtet. In sämtlichen Arbeiten, die eine solche Ausrichtung vornehmen, kommt dafür \textit{Dynamic Time Warping} zum Einsatz \cite{Tharatipyakul2024}.
Die vorliegende Arbeit folgt dem regelbasierten Ansatz. Ausschlaggebend sind zwei Gründe. Erstens ist die Nachvollziehbarkeit der Rückmeldung das zentrale Ziel der Anwendung, da eine Bewertung ohne erkennbare Begründung im Trainingskontext wenig Nutzen bringt. Zweitens stand kein annotierter Datensatz zur Verfügung, der die erforderlichen Angaben zur Ausführung und Qualität der betrachteten Taekwondo-Techniken in geeigneter Form bereitstellt. Die daraus abgeleiteten Kriterien und Schwellenwerte werden in Kapitel~4 hergeleitet.
\section{Verwandte Arbeiten}
Systeme zur automatisierten Bewertung von Bewegungsausführungen sind für verschiedene Anwendungsfelder entwickelt worden. Sie unterscheiden sich darin, welche Sensorik sie voraussetzen, ob sie die ausgeführte Bewegung selbst erkennen müssen und woher der Maßstab stammt, an dem sie die Ausführung messen.
Pham et al. \cite{Pham2022} stellten eine Anwendung vor, die körperliche Übungen aus RGB-Bildern erkennt und bewertet. Das System gewinnt die \textit{Landmarks} mit MediaPipe, ordnet die Sequenz über ein neuronales Netz einer von drei Übungen zu und bewertet sie anschließend. Die Arbeit zeigt, dass \textit{Pose Estimation}, Erkennung der Übung und die Bewertung ohne spezielle Hardware auskommen und damit auf Geräte wie die Kinect verzichten können, die vergleichbare Systeme zuvor voraussetzten \cite{Pham2022}.
Wie belastbar die dabei erzeugten Bewertungen sind, ist dagegen wenig untersucht. Ein systematisches Review kamerabasierter mobiler Anwendungen für das Bewegungsscreening bei gesunden Erwachsenen weist darauf hin, dass zur Bewertung der Bewegungsqualität und zu Rückmeldungen in Echtzeit bislang wenig Evidenz vorliegt \cite{ElRajab2025}. Die technische Machbarkeit ist damit vielfach gezeigt, die Frage nach der Aussagekraft der Ergebnisse bleibt weitgehend offen.
Wie in Abschnitt~2.1 dargestellt, verteilt sich die biomechanische Forschung zu Taekwondo-Kicks ungleichmäßig über die Techniken. Gerade die Technik mit der stärksten Körperdrehung ist dabei am schlechtesten untersucht. Kim et al. \cite{Kim2011} untersuchten die Gelenkkinematik des Kickbeins für vier Taekwondo-Techniken und beschreiben den \textit{Spinning Back Kick} als \textit{push-like movement}, das seine Geschwindigkeit aus einer Kombination von Hüft- und Kniestreckung gewinnt \cite{Kim2011}. Auf diesen Befund stützt sich die Herleitung der Bewertungskriterien für den \textit{Spinning Back Kick} in Abschnitt~4.8.
\subsection{Wang et al. (2026)}
Für Taekwondo liegt seit Kurzem eine Arbeit vor, die dem hier verfolgten Ansatz nahekommt. Wang et al. \cite{Wang2026} entwickelten einen Desktop-Prototyp, der mit dem Smartphone aufgenommene Videos auswertet und drei Grundtechniken erkennt und bewertet, nämlich \textit{Front Kick}, \textit{Roundhouse Kick} und \textit{Axe Kick}. Grundlage ist ein eigens erhobener Datensatz mit 765 verwertbaren Aufnahmen von 150 Studierenden \cite{Wang2026}. Die Verarbeitungskette ist in Abbildung \ref{fig:wang_pipeline} dargestellt.
\begin{figure}[htbp]
    \centering
    \includegraphics[width=\textwidth]{Images/wang\_pipeline.png}
    \caption[Verarbeitungskette bei Wang et al.]{%
    Verarbeitungskette des Prototyps von Wang et al.: Merkmalsgewinnung,
    Technikerkennung, regelbasierte Bewertung und Rückmeldung. Abbildung
    entnommen aus \cite{Wang2026}}
    \label{fig:wang_pipeline}
\end{figure}
Das System verwendet MediaPipe Pose zur Extraktion von 13 Landmarken, aus denen je Bild ein 52-dimensionaler Merkmalsvektor aus normierten zweidimensionalen Positionen und deren erster Ableitung gebildet wird. Die Erkennung der Technik übernimmt ein bidirektionales LSTM, das unter dem verwendeten Validierungsverfahren eine Genauigkeit von 75,3 $\pm$ 2,3 Prozent erreicht. Die Autoren ordnen diesen Wert selbst vorsichtig ein, da in den Daten keine Teilnehmerkennungen vorlagen und sich deshalb keine personenunabhängige Trennung von Trainings- und Validierungsdaten herstellen ließ. Aufnahmen derselben Person können daher in beiden Teilmengen vorkommen, was die Genauigkeit überhöht haben kann. Transformer- und GCN-Architekturen wurden zum Vergleich untersucht und schnitten mit 67,3 beziehungsweise 51,0 Prozent schlechter ab \cite{Wang2026}.
Aufschlussreich für die vorliegende Arbeit ist die Aufteilung über die drei Techniken. Der \textit{Roundhouse Kick} wird mit einem F1-Wert von 67,9 Prozent deutlich schlechter erkannt als der \textit{Front Kick} mit 73,9 und der \textit{Axe Kick} mit 81,0 Prozent. Als Ursache nennen die Autoren die stärkere Rotation bei dieser Technik, die die Abhängigkeit vom Blickwinkel erhöht und zur Selbstverdeckung von Hüfte, Knie und Sprunggelenk führt \cite{Wang2026}. Einschränkend kommt hinzu, dass der \textit{Roundhouse Kick} zugleich die kleinste Klasse bildet und ausschließlich von Studierenden des zweiten Jahrgangs stammt, da die Technik im ersten Jahrgang noch nicht im Unterricht behandelt worden war \cite{Wang2026}. 
Die Bewertung selbst erfolgt regelbasiert nach einem 10-Punkte-Schema, das an die Bewertungsmaßstäbe des Poomsae-Wettkampfs angelehnt ist und sich in vier Punkte für technische Korrektheit sowie sechs Punkte für Ausdruck aufteilt \cite{Wang2026}. Die Ausdruckskomponente stützt sich auf abgeleitete Größen wie Geschwindigkeit und Gleichmäßigkeit der Kicks und bleibt hier außer Betracht, da die vorliegende Arbeit auf eine solche Komponente verzichtet. 
Die Genauigkeitskomponente umfasst sechs Positionen: Streckung des Standbeins, Kicktechnik aus erreichter Höhe und Kniestreckung, Chamber, Rückzug, Landeposition und Stabilität der Hände \cite{Wang2026}. Vier davon finden sich auch in der vorliegenden Arbeit wieder. Landeposition und Handhaltung fehlen, da weder die \textit{Landmarks} am Fuß noch am Handgelenk ausgewertet werden. Beide werden in Kapitel~6 als mögliche Erweiterung aufgegriffen.

\subsection{Abgrenzung der vorliegenden Arbeit}
Aus den vorgestellten Arbeiten lässt sich die Position der vorliegenden Arbeit ableiten. Pham et al. \cite{Pham2022} zeigen, dass eine markerlose, MediaPipe-gestützte Bewertung ohne Spezialhardware grundsätzlich funktioniert. Ihr System zielt allerdings auf zyklische Fitnessübungen und nicht auf schnelle, rotationsbehaftete Kampfsporttechniken. El-Rajab et al. \cite{ElRajab2025} machen deutlich, dass die Aussagekraft solcher Bewertungen selten eigenständig untersucht wird. Dieser Punkt gilt auch für die vorliegende Arbeit als Grenze, denn eine empirische Kalibrierung an Trainerurteilen wäre dafür erforderlich und konnte hier nicht geleistet werden. Kapitel~6 greift ihn wieder auf.
Am nächsten steht der vorliegenden Arbeit der Ansatz von Wang et al. \cite{Wang2026}. Zwei Unterschiede sind für die Zielsetzung dieser Arbeit zentral.
Erstens verzichtet die vorliegende Arbeit auf die automatische Technikerkennung und lässt die Technik stattdessen vom Nutzer wählen. Damit entfällt die Fehlerquelle einer LSTM-basierten Klassifikation, die bei Wang et al. gerade für die am stärksten rotierende Technik am unzuverlässigsten arbeitet. 
Zweitens wird mit dem \textit{Spinning Back Kick} eine Technik bewertet, deren Rotation deutlich stärker ausfällt als bei den drei von Wang et al. untersuchten Techniken. Weder der \textit{Front Kick} noch der \textit{Axe Kick} besitzen eine vergleichbare Drehung um die Körperlängsachse. Auch der bei Wang et al. am schlechtesten erkannte \textit{Roundhouse Kick} erfordert keine etwa 180 Grad umfassende Wendung des gesamten Körpers. Der \textit{Spinning Back Kick} ermöglicht es daher, die Grenzen des gewählten regelbasierten Ansatzes an einer stärker rotationsgeprägten Technik zu untersuchen.

Daraus ergibt sich die Abgrenzung der vorliegenden Arbeit. Untersucht wird, wie ausgewählte Merkmale der drei Techniken aus acht \textit{Landmarks} abgeleitet und regelbasiert bewertet werden können und welche Grenzen sich dabei insbesondere bei einer ausgeprägten Körperdrehung von etwa 180 Grad zeigen. Die Arbeit beansprucht dabei keine unabhängige Validierung der Messgenauigkeit, sondern untersucht die prototypische Umsetzbarkeit und die daraus hervorgehenden Grenzen.
\section{Konzept \& Implementierung}
\subsection{Überblick über das System}
Die entwickelte Anwendung nimmt ein Video einer einzelnen Kickausführung entgegen und gibt eine nach Kriterien aufgeschlüsselte Bewertung zurück. Der Ablauf gliedert sich in vier Schritte.
Im ersten Schritt wird die Aufnahme Bild für Bild durchlaufen und für jedes Einzelbild die Position der ausgewerteten \textit{Landmarks} bestimmt. MediaPipe arbeitet dabei intern in zwei Stufen, indem zunächst ein Detektor die Person im Bild findet und anschließend ein Tracker die \textit{Landmarks} innerhalb dieses Bildausschnitts lokalisiert (Abschnitt~2.3.3).
Im zweiten Schritt werden daraus kinematische Größen berechnet, insbesondere Gelenkwinkel und Positionen einzelner Körperpunkte. Der dritte Schritt teilt die Aufnahme in die Phasen des jeweiligen Bewegungsablaufs, sodass jede Messgröße dem Abschnitt zugeordnet werden kann, in dem sie aussagekräftig ist. Im vierten Schritt werden die gewonnenen Werte mit festgelegten Schwellen verglichen und in eine Bewertung überführt.
Umgesetzt ist das System als Webanwendung mit getrenntem Frontend und Backend. Das Frontend übernimmt die Auswahl von Technik und Videodatei, die Anzeige des Verarbeitungsfortschritts sowie die Darstellung der Ergebnisse. Die Analyse erfolgt serverseitig in Python. Im Betrieb wird das Frontend vom Backend statisch mit ausgeliefert, sodass Anwendung und Schnittstelle unter derselben Adresse erreichbar sind.
Diese Aufteilung wurde bewusst gewählt. Eine Auswertung im Browser oder über eine App auf dem Smartphone wäre technisch möglich gewesen, hätte die verfügbare Rechenleistung jedoch an das jeweilige Endgerät gebunden und den Einsatz der größten Modellvariante von MediaPipe praktisch ausgeschlossen. Die serverseitige Verarbeitung stellt zudem sicher, dass identische Aufnahmen unabhängig vom verwendeten Gerät gleich ausgewertet werden.
Abbildung \ref{fig:architecture_diagram} zeigt die Architektur des Systems. Sie besteht aus drei Teilen, dem React-basierten Frontend, der FastAPI-Web-API und der serverseitigen Python-Analysepipeline. Das Frontend übermittelt die ausgewählte Videodatei und die gewählte Kicktechnik an das Backend. Dieses startet die Analyse und übergibt das Video an die Analysepipeline, in der \textit{Pose Estimation}, Phasenerkennung und Bewertung ablaufen. Anschließend werden die Ergebnisse und das annotierte Video über das Backend an das Frontend zurückgegeben und dort angezeigt.
\begin{figure}[h]
    \centering
    \includegraphics[width=1\textwidth]{Images/diagramm.png}
    \caption{Vereinfachte Architektur der entwickelten Anwendung.}
    \label{fig:architecture_diagram}
\end{figure}
\subsubsection{Technologische Grundlage}
Das Backend ist in Python umgesetzt und verwendet FastAPI als Webframework sowie Uvicorn als Server. Die Bildverarbeitung erfolgt mit OpenCV, während numerische Berechnungen mit NumPy durchgeführt werden. Für die \textit{Pose Estimation} wird MediaPipe eingesetzt. Die abschließende Kodierung des Ausgabevideos erfolgt mit \texttt{imageio-ffmpeg}. Das Frontend ist mit React und TypeScript umgesetzt, als Buildwerkzeug wird Vite verwendet. Die Anwendung ist als Progressive Web App ausgelegt und kann dadurch sowohl auf Desktop-Systemen im Browser als auch auf mobilen Endgeräten genutzt werden.

Für die Umsetzung wurde Python 3.14 verwendet. Die während der Entwicklung eingesetzte MediaPipe-Version war für diese Python-Version nicht offiziell ausgewiesen. Der Betrieb funktionierte in der vorliegenden Entwicklungsumgebung dennoch. Für eine reproduzierbare Weiterentwicklung sollte die verwendete Kombination aus Python und MediaPipe deshalb über feste Versionsangaben dokumentiert und in einer offiziell unterstützten Umgebung erneut geprüft werden.

\subsection{Erfassung der Körperhaltung}
\subsubsection{Wahl des Frameworks}
Für die Erfassung der Körperhaltung wird MediaPipe eingesetzt. Ausschlaggebend waren drei Eigenschaften. Das Verfahren ist für die Verarbeitung auf handelsüblicher Hardware ausgelegt und liefert neben normierten Bildkoordinaten auch dreidimensionale Weltkoordinaten, aus denen sich Gelenkwinkel ableiten lassen. Zudem wird MediaPipe unter der Apache License 2.0 veröffentlicht und erlaubt damit auch eine weitergehende Nutzung unter Einhaltung der Lizenzbedingungen \cite{MediaPipeLicense}. OpenPose wird dagegen unter einer Lizenz der Carnegie Mellon University bereitgestellt, die die Nutzung auf nichtkommerzielle Forschung an akademischen oder gemeinnützigen Einrichtungen begrenzt \cite{OpenPoseLicense}.

Eingesetzt wird der \textit{Pose Landmarker} im Laufmodus \texttt{VIDEO}, der die in Abschnitt~2.3.3 beschriebene Verfolgung über die Bildfolge hinweg nutzt. Die Konfidenzschwellen für Detektion, Präsenz und Tracking bleiben auf dem voreingestellten Wert von 0,5.

Alle drei Modellvarianten sind im Prototyp auswählbar. Voreingestellt ist \textit{Heavy}, da diese Variante in Vorversuchen die subjektiv stabilsten Landmark-Verläufe lieferte. Diese Beobachtung stellt keine systematische Modellvalidierung dar. Die Auswahlmöglichkeit dient deshalb in erster Linie Entwicklungs und Testzwecken. Für eine Endanwendung wäre eine solche Modellauswahl nicht erforderlich. Ein kontrollierter Vergleich derselben Videos mit \textit{Lite}, \textit{Full} und \textit{Heavy} wird in Abschnitt~\ref{sec:ausblick} als möglicher Validierungsschritt beschrieben.

\subsubsection{Verwendete Landmarks}
Von den 33 gelieferten \textit{Landmarks} werden acht ausgewertet, nämlich Schulter, Hüfte, Knie und Sprunggelenk auf beiden Körperseiten. Diese Punkte ermöglichen die Berechnung der in der vorliegenden Arbeit verwendeten Kniewinkel, Hüftwinkel und Richtungsgrößen. Der Kniewinkel ergibt sich aus Hüfte, Knie und Sprunggelenk, die Hüftbeugung aus Schulter, Hüfte und Knie. Für die Ausrichtung des Beckens wird die Verbindungslinie beider Hüften verwendet. Aus ihrer Senkrechten wird in der Draufsicht eine körperrelative Beckenrichtung abgeleitet (Abschnitt~4.7).

Die Landmarken an Ferse und Zehen bleiben unberücksichtigt. Dadurch können weder die Ausrichtung des kickenden Fußes noch der Pivot des Standfußes direkt erfasst werden. Welche zusätzlichen Kriterien dadurch entfallen, wird in Kapitel~6 diskutiert.

Sämtliche für die Bewertung verwendeten Winkel und Positionen stützen sich auf die dreidimensionalen Weltkoordinaten. Die normierten Bildkoordinaten werden ausschließlich für die Darstellung des Skeletts im Ausgabevideo verwendet. MediaPipe liefert zusätzlich Sichtbarkeitswerte für die einzelnen Landmarken. Diese Werte werden in der aktuellen Implementierung nicht zur Gewichtung oder Verwerfung von Messpunkten genutzt und stellen damit eine weitere mögliche Robustheitsverbesserung dar.

\subsection{Vorverarbeitung}
Zwischen der hochgeladenen Datei und den ausgewerteten Messgrößen liegen mehrere Verarbeitungsschritte, die das Ergebnis beeinflussen.
\subsubsection{Bildrate}
Das Video wird mit OpenCV geöffnet und die Bildrate aus den Metadaten ausgelesen. Liefert die Datei keine verwertbare Angabe, wird ersatzweise mit 30 Bildern pro Sekunde gerechnet. Die Bildrate geht ausschließlich in die Berechnung der Fußgeschwindigkeit ein (Abschnitt~4.11), während die winkelbasierten Kriterien von ihr unabhängig sind. Ein falscher Ersatzwert verfälscht damit nur die angezeigte Geschwindigkeit, nicht aber die Bewertung. 
\subsubsection{Auswahl der Person}
Werden mehrere Personen im Bild erkannt, wird stets die erste zurückgegebene Person verwendet. Eine Auswahl durch den Nutzer oder ein Hinweis findet nicht statt. Dass sich nur eine Person im Bild befinden soll, wird ausschließlich als Hinweistext vermittelt und technisch nicht geprüft (Abschnitt~4.14).
Das System prüft nicht, ob über das ganze Video hinweg dieselbe Person ausgewertet wird. Bricht die Verfolgung kurz ab, läuft die Detektion erneut an und kann eine andere Person zurückgeben (Abschnitt~2.3.3). Da immer die erste zurückgegebene Person verwendet wird, kann die Auswertung dadurch mitten in der Aufnahme auf jemand anderen überspringen.
\subsubsection{Normierung}
Die Weltkoordinaten von MediaPipe sind bereits auf die Hüftmitte bezogen und damit translationsnormiert (Abschnitt~2.3.3). Eine Skalierungsnormierung ist für die Gelenkwinkel nicht erforderlich, da diese skaleninvariant sind. Der Winkel zwischen zwei Segmenten ändert sich nicht, wenn beide Segmente länger werden.
Für längenbasierte Größen gilt das nicht. Die Fußbahn, die beim \textit{Spinning Back Kick} zur Phasenerkennung dient, wird deshalb auf die Beinlänge normiert, die sich als Summe der Abstände Hüfte-Knie und Knie-Knöchel des Standbeins ergibt. 
\begin{equation}
\ell_{\text{Bein}} = \lVert \vec{p}_{\text{Hüfte}} - \vec{p}_{\text{Knie}} \rVert + \lVert \vec{p}_{\text{Knie}} - \vec{p}_{\text{Knöchel}} \rVert
\label{eq:leglength}
\end{equation}
Dabei beziehen sich die Punkte $\vec{p}_{\text{Hüfte}}$, $\vec{p}_{\text{Knie}}$ und $\vec{p}_{\text{Knöchel}}$ auf das Standbein und gehen als dreidimensionale Weltkoordinaten in Metern ein. Alle Größen der Fußbahn werden anschließend in Vielfachen von $\ell_{\text{Bein}}$ angegeben (Abschnitt~4.4).  Eine Fußhöhe von $0{,}30\,\ell_{\text{Bein}}$ bedeutet damit, dass der Fuß auf 30 Prozent der Beinlänge angehoben ist, unabhängig davon, wie lang das Bein tatsächlich ist. Die Kriterien zur Fußbahn sind dadurch von der Körpergröße unabhängig.
Die Kniehöhe wird dagegen in absoluten Metern ausgewertet. Bewertet wird das Absinken des Knies während der Kick-Phase, also die Differenz zwischen der höchsten und der niedrigsten erreichten Position. 
\subsubsection{Glättung}
Nicht alle Messgrößen werden geglättet. Die Fußgeschwindigkeit wird über einen gleitenden Median von drei Bildern geführt, da sie als Differenzgröße besonders rauschempfindlich ist. Die übrigen Winkel gehen ungeglättet in die Messreihe ein. Beim Kniewinkel des Standbeins wird die Streuung erst bei der Bewertung aufgefangen, indem dort der Median über alle Bilder der Kick-Phase gebildet wird.
Zusätzlich wirken zwei Plausibilitätsfilter. Sie glätten nicht, sondern verwerfen Werte, die keiner echten Bewegung entsprechen können. MediaPipe vertauscht gelegentlich linke und rechte \textit{Landmarks}, wodurch die berechnete Beckenrichtung schlagartig um 180 Grad kippt. 
Der erste Filter betrifft die Hüftausrichtung. Ändert sie sich zwischen zwei aufeinanderfolgenden Bildern um mehr als 60 Grad, wird der neue Wert verworfen und der letzte gültige beibehalten. Bei 30 Bildern pro Sekunde entspricht diese Grenze einer Winkelgeschwindigkeit von 1800 Grad pro Sekunde, die eine echte Beckenbewegung nicht erreicht. Der Filter wirkt bei allen drei Techniken, da alle die Hüftausrichtung im vollen Bereich von 0 bis 180 Grad auswerten (Abschnitt~4.7).
Der zweite Filter betrifft die aufsummierte Beckendrehung und wirkt allein beim \textit{Spinning Back Kick}, da nur dort eine solche Summe gebildet wird. Hier liegt die Grenze enger bei 30 Grad je Bild. Übersteigt ein Schritt diesen Wert, wird er nicht mitgezählt. Weil sich die Einzelschritte aufsummieren, würde ein einzelner fehlerhafter Wert sonst dauerhaft im Ergebnis stehen bleiben.
\subsection{Erkennung der Bewegungsphasen}
\subsubsection{Warum eine Phaseneinteilung erforderlich ist}
Ob ein Messwert für eine gute oder eine schlechte Ausführung spricht, hängt davon ab, wann im Bewegungsablauf er auftritt. Ein Kniewinkel von etwa 70 Grad ist in der Chamber-Phase erwünscht, denn dort soll das Bein eng angezogen sein. Derselbe Wert während der Streckung würde bedeuten, dass der Kick gar nicht ausgeführt wurde. Umgekehrt ist ein weit gestrecktes Knie in der Kick-Phase das Ziel und in der Chamber-Phase ein Zeichen dafür, dass nicht angezogen wurde.
Die Messgrößen lassen sich deshalb nur bewerten, wenn bekannt ist, zu welchem Abschnitt der Bewegung sie gehören. Die Phaseneinteilung ist damit die Voraussetzung dafür, dass die Kriterien überhaupt gebildet werden können.
\subsubsection{Umsetzung als Zustandsmaschine}
Jede Technik ist als eigene Zustandsmaschine umgesetzt. Der Zustand beschreibt die aktuelle Phase, die Übergänge sind über Bedingungen auf den Messgrößen des laufenden Bildes definiert (Abschnitt~4.4.4). Jeder Phasenwechsel wird mit der zugehörigen Bildnummer protokolliert, sodass sich der erkannte Bewegungsablauf im Nachhinein nachvollziehen lässt.
\textit{Front Kick} und \textit{Roundhouse Kick} durchlaufen die Phasenfolge \texttt{idle}, \texttt{chamber}, \texttt{kick}, \texttt{rechamber} und zurück nach \texttt{idle}. Der \textit{Spinning Back Kick} besitzt mit \texttt{rotation} eine zusätzliche Phase, in der die Eindrehung des Körpers stattfindet, bevor das Knie angezogen wird.
\subsubsection{Bestimmung der Kickseite}
Vor der eigentlichen Phasenerkennung muss festgelegt werden, welches Bein kickt und welches steht. Solange sich das System in der Phase \texttt{idle} befindet, vergleicht das System die Höhe beider Knie. Übersteigt sie fünf Zentimeter, gilt das höhere Knie als Kickbein. Diese Zuordnung wird danach nicht mehr zurückgesetzt und bleibt für das gesamte Video bestehen.
Damit lässt sich pro Video nur eine Ausführung auswerten. Die Beschränkung folgt aus der Zielsetzung. Untersucht werden soll, wie zuverlässig sich die Merkmale einer einzelnen Ausführung erfassen lassen. Eine kontrollierte Einzelaufnahme ist dafür die geeignete Grundlage, während mehrere Kicks in einer Aufnahme eine weitere Fehlerquelle einführen würden. 
\subsubsection{Übergangsbedingungen}
Die Übergänge zwischen den Phasen werden für jede Technik über Bedingungen auf den Messgrößen des jeweils aktuellen Bildes definiert. Die verwendeten Schwellen dienen der Zustandszuordnung und sind von den späteren Bewertungsschwellen zu unterscheiden. Ein Teil der Werte orientiert sich an den Gelenkwinkeln, die auch in der Bewertung verwendet werden. Die übrigen Übergangsschwellen wurden im Rahmen der prototypischen Entwicklung heuristisch anhand der beschriebenen Bewegungsstruktur und eigener Testaufnahmen festgelegt. Sie wurden nicht anhand eines annotierten Datensatzes empirisch optimiert. Diese Abhängigkeit von heuristischen Übergängen ist eine konzeptionelle Grenze des Systems und wird in Kapitel~6 erneut aufgegriffen.

\paragraph{Front Kick.}
Die Chamber-Phase beginnt, sobald der Kniewinkel unter 100 Grad fällt und der Oberschenkel gleichzeitig um mehr als 45 Grad gegenüber der Senkrechten angehoben ist. Die zweite Bedingung verhindert, dass ein bloßes Beugen des Knies im Stand als Chamber erkannt wird. Der Übergang in die Kick-Phase erfolgt, wenn der Kniewinkel 131 Grad überschreitet und die Oberschenkelneigung weiterhin über 45 Grad liegt. Der Rückzug gilt als eingeleitet, sobald der Kniewinkel um mindestens 80 Grad unter das erreichte Streckungsmaximum zurückgeht, während die Oberschenkelneigung weiterhin über 45 Grad bleibt. Fällt die Oberschenkelneigung unter 50 Grad, während das Knie über 150 Grad gestreckt ist, wird die Erkennung abgebrochen und in den Ausgangszustand zurückgesetzt. Tabelle~\ref{tab:transitions_front} fasst die Übergänge zusammen.

\begin{table}[htbp]
\centering
\begin{tabular}{p{0.32\textwidth} p{0.62\textwidth}}
\hline
\textbf{Übergang} & \textbf{Bedingung} \\ \hline
\texttt{idle} $\rightarrow$ \texttt{chamber} & Kniewinkel $<100°$ \textbf{und} Oberschenkelneigung $>45°$ \\
\texttt{chamber} $\rightarrow$ \texttt{kick} & Kniewinkel $>131°$ \textbf{und} Oberschenkelneigung $>45°$ \\
\texttt{kick} $\rightarrow$ \texttt{rechamber} & Oberschenkelneigung $>45°$ \textbf{und} (max. Kniewinkel $-$ aktuell) $\geq80°$ \\
$\rightarrow$ \texttt{idle} (Abbruch) & Kniewinkel $>150°$ \textbf{und} Oberschenkelneigung $<50°$ \\ \hline
\end{tabular}
\caption{Phasenübergänge des \textit{Front Kicks}}
\label{tab:transitions_front}
\end{table}

Die Schwelle von 131 Grad für den Übergang in die Kick-Phase entspricht dem \texttt{fail\_at} Wert des Kriteriums Beinstreckung (Abschnitt~4.8). Eine gerade noch erkannte Ausführung beginnt damit bei 0 Prozent für dieses Kriterium. Die fachliche Herleitung der Bewertungsschwellen erfolgt getrennt von der Phasenerkennung in Abschnitt~4.8.

\paragraph{Roundhouse Kick.}
Der \textit{Roundhouse Kick} folgt derselben Zustandsstruktur wie der \textit{Front Kick} und verwendet für Chamber und Kick dieselben Schwellen. Die Phasenerkennung stützt sich dabei nur auf Kniewinkel und Oberschenkelneigung. Die technikspezifische Beckenrotation wird nicht zur Erkennung der Phase herangezogen, sondern erst über die körperrelative Hüftausrichtung in der Bewertung berücksichtigt (Abschnitt~4.8).

\begin{table}[htbp]
\centering
\begin{tabular}{p{0.32\textwidth} p{0.62\textwidth}}
\hline
\textbf{Übergang} & \textbf{Bedingung} \\ \hline
\texttt{idle} $\rightarrow$ \texttt{chamber} & Kniewinkel $<100°$ \textbf{und} Oberschenkelneigung $>45°$ \\
\texttt{chamber} $\rightarrow$ \texttt{kick} & Kniewinkel $>131°$ \textbf{und} Oberschenkelneigung $>45°$ \\
\texttt{kick} $\rightarrow$ \texttt{rechamber} & Oberschenkelneigung $>45°$ \textbf{und} (max. Kniewinkel $-$ aktuell) $\geq80°$ \\
$\rightarrow$ \texttt{idle} (Abbruch) & Kniewinkel $>150°$ \textbf{und} Oberschenkelneigung $<50°$ \\ \hline
\end{tabular}
\caption{Phasenübergänge des \textit{Roundhouse Kicks}}
\label{tab:transitions_round}
\end{table}

\paragraph{Spinning Back Kick.}
Der \textit{Spinning Back Kick} besitzt als einzige der drei Techniken eine vorgelagerte Rotationsphase. Sie beginnt, sobald die aufsummierte Änderung der Beckenrichtung 45 Grad überschreitet. Bleibt dieser Wert darunter, wird das Knie aber dennoch deutlich angezogen, greift ein Rückfall direkt von \texttt{idle} nach \texttt{chamber}. Dadurch kann auch eine Ausführung analysiert werden, bei der die vorgelagerte Rotationsphase nicht zuverlässig erkannt wurde.

Anders als bei den beiden übrigen Techniken wird der Übergang in die Kick-Phase nicht über den Kniewinkel erkannt, sondern über die Fußbahn. Der horizontale Abstand der beiden Sprunggelenke muss über $0{,}45\,\ell_{\text{Bein}}$ liegen und die Fußhöhe über $0{,}30\,\ell_{\text{Bein}}$. Der Fuß muss damit gleichzeitig deutlich vom Standbein entfernt und angehoben sein.

Der Kniewinkel allein ist für diesen Übergang nicht eindeutig genug. Eine weitgehende Streckung kann sowohl während des Kicks als auch beim späteren Absetzen auftreten. Die zusätzliche Berücksichtigung der Fußposition soll diese Zustände voneinander trennen.

\begin{table}[htbp]
\centering
\begin{tabular}{p{0.30\textwidth} p{0.58\textwidth}}
\hline
\textbf{Übergang} & \textbf{Bedingung} \\ \hline
\texttt{idle} $\rightarrow$ \texttt{rotation} & aufsummierte Änderung der Beckenrichtung $>45°$ \\
\texttt{idle} $\rightarrow$ \texttt{chamber} (Rückfall) & Kniewinkel $<100°$ \textbf{und} Hüftbeugung $<150°$ \\
\texttt{rotation} $\rightarrow$ \texttt{chamber} & Kniewinkel $<100°$ \\
\texttt{chamber} $\rightarrow$ \texttt{kick} & Sprunggelenkabstand $>0{,}45\,\ell_{\text{Bein}}$ \textbf{und} Fußhöhe $>0{,}30\,\ell_{\text{Bein}}$ \\
\texttt{chamber} $\rightarrow$ \texttt{idle} (Abbruch) & Fuß war oben ($>0{,}30\,\ell_{\text{Bein}}$), jetzt unten ($<0{,}15\,\ell_{\text{Bein}}$) \\
\texttt{kick} $\rightarrow$ \texttt{rechamber} & Fußhöhe $>0{,}25\,\ell_{\text{Bein}}$ \textbf{und} (max. Kniewinkel $-$ aktuell) $\geq35°$ \\
\texttt{kick} $\rightarrow$ \texttt{idle} & Kniewinkel $>160°$ \textbf{und} Hüftbeugung $>160°$ \\
\texttt{rechamber} $\rightarrow$ \texttt{idle} & Kniewinkel $>160°$ \\ \hline
\end{tabular}
\caption{Phasenübergänge des \textit{Spinning Back Kicks}. $\ell_{\text{Bein}}$ bezeichnet die Beinlänge aus Gleichung~\ref{eq:leglength}.}
\label{tab:transitions_dwit}
\end{table}

Nur beim \textit{Spinning Back Kick} wird zusätzlich eine aufsummierte Änderung der Beckenrichtung geführt. Der in Abschnitt~4.3.4 beschriebene Plausibilitätsfilter verwirft dabei einzelne Änderungen von mehr als 30 Grad pro Bild. Da diese Größe die globale Körperdrehung nur näherungsweise abbildet, wird sie nicht direkt als Bewertungskriterium verwendet.

\subsection{Treffmoment und technikspezifische Sonderregeln}
Alle drei Techniken bestimmen einen Treffmoment auf dieselbe Weise. Während der Kick-Phase speichert das System den größten bisher gemessenen Kniewinkel. Wird ein neuer Höchstwert erreicht, wird zusätzlich die aktuelle Ausrichtung des Kickbeins relativ zum Becken gespeichert. Am Ende der Kick-Phase entspricht dieser Wert damit der Ausrichtung in dem Bild mit der stärksten gemessenen Beinstreckung. Er wird anschließend in die technikspezifische Bewertung einbezogen (Abschnitt~4.8).

Dieser Zeitpunkt wird in der vorliegenden Arbeit als Treffmoment angenähert. Die Annahme lautet, dass die stärkste Beinstreckung bei den betrachteten Techniken in der Nähe des Zeitpunkts liegt, an dem ein Treffer stattfinden würde. Da die Videos ohne tatsächliches Ziel aufgenommen werden, kann der reale Kontaktzeitpunkt nicht beobachtet werden. Der Begriff Treffmoment bezeichnet daher keinen gemessenen physischen Kontakt, sondern den für die Bewertung gewählten Referenzframe.

Die drei Techniken unterscheiden sich darin, welche Ausrichtung des Kickbeins relativ zum Becken in diesem Referenzframe angestrebt wird (Tabelle~\ref{tab:hip_alignment_technique}). Beim \textit{Front Kick} soll das Kickbein annähernd in Beckenblickrichtung geführt werden, weshalb ein kleiner Winkel angestrebt wird. Beim \textit{Roundhouse Kick} wird eine seitliche Ausrichtung um 90 Grad angestrebt. Beim \textit{Spinning Back Kick} soll das Kickbein im Verhältnis zum Becken nach hinten ausgerichtet sein, sodass der Winkel in Richtung 180 Grad zunimmt.

Beim \textit{Roundhouse Kick} liegt der Zielwert mit 90 Grad in der Mitte des Wertebereichs. Deshalb wird nicht der Rohwert direkt bewertet, sondern dessen absolute Abweichung von 90 Grad (Gleichung~\ref{eq:rotationerror}). Beim \textit{Front Kick} und beim \textit{Spinning Back Kick} kann der Rohwert dagegen unmittelbar in die monotone Bewertungsfunktion eingehen.

Die Verwendung eines einzelnen Referenzframes verhindert, dass eine spätere Weiterdrehung nach der stärksten Beinstreckung die Ausrichtung im vermeintlichen Treffzeitpunkt ersetzt. Gleichzeitig erhöht sie die Empfindlichkeit gegenüber Ausreißern in genau diesem Bild. Der Treffmoment stellt damit einen Kompromiss zwischen phasenspezifischer Bewertung und Robustheit gegenüber einzelnen fehlerhaften Landmark-Schätzungen dar. Diese Grenze wird in Kapitel~6 erneut aufgegriffen.

\subsection{Bewertungsfunktion}
Alle Bewertungskriterien werden mit derselben linearen Funktion auf einen Wert zwischen 0 und 100 Prozent abgebildet:

\begin{equation}
\text{score}(v) = 100 \cdot \max\!\left(0,\; \min\!\left(1,\; \frac{v - v_{\text{fail}}}{v_{\text{ideal}} - v_{\text{fail}}}\right)\right)
\label{eq:score}
\end{equation}

Dabei bezeichnet $v$ den gemessenen Wert. $v_{\text{fail}}$ markiert den Rand des Bereichs, an dem das Kriterium 0 Prozent erhält. $v_{\text{ideal}}$ bezeichnet den Wert, ab dem 100 Prozent erreicht werden. Die äußeren Funktionen begrenzen die Ausgabe auf den Bereich zwischen 0 und 100 Prozent.

Die Formel kann sowohl für Kriterien verwendet werden, bei denen ein größerer Wert besser ist, als auch für Kriterien, bei denen ein kleinerer Wert besser ist. Bei der Beinstreckung gilt beispielsweise $v_{\text{fail}}=131$ Grad und $v_{\text{ideal}}=170$ Grad. Beim Chamber-Winkel ist ein kleinerer Wert erwünscht. Hier gilt $v_{\text{fail}}=100$ Grad und $v_{\text{ideal}}=71$ Grad. Ein gemessener Chamber-Winkel von 85 Grad ergibt damit rund 52 Prozent.

Ein Kriterium wird in der Benutzeroberfläche ab 50 Prozent als bestanden behandelt. Diese Bestehensgrenze ist von $v_{\text{fail}}$ und $v_{\text{ideal}}$ zu unterscheiden. Sie beeinflusst lediglich, ob ein Verbesserungshinweis angezeigt wird. Die berechnete Punktzahl selbst bleibt unverändert.

Die Bewertungsfunktion legt damit nur fest, wie ein Messwert zwischen zwei Grenzwerten in einen Prozentwert überführt wird. Ob diese Grenzwerte die technische Qualität eines Kicks fachlich angemessen abbilden, ist eine davon getrennte Frage. Die Grenzwerte wurden in der vorliegenden Arbeit nicht anhand unabhängiger Trainerurteile empirisch kalibriert.

\subsection{Winkelberechnung}
Alle Gelenkwinkel werden dreidimensional aus den Weltkoordinaten berechnet. Ein Gelenkwinkel wird durch drei \textit{Landmarks} beschrieben, nämlich zwei äußere Punkte $\vec{a}$ und $\vec{c}$ sowie den Scheitelpunkt $\vec{b}$. Beim Kniewinkel sind dies Hüfte, Knie und Sprunggelenk, wobei das Knie den Scheitelpunkt bildet. Von diesem Punkt aus werden die Vektoren $\vec{u}=\vec{a}-\vec{b}$ und $\vec{w}=\vec{c}-\vec{b}$ gebildet. Der eingeschlossene Winkel $\theta$ ergibt sich über das normierte Skalarprodukt:

\begin{equation}
\theta = \arccos\!\left(\frac{\vec{u}\cdot\vec{w}}{\lVert\vec{u}\rVert\,\lVert\vec{w}\rVert}\right)
\label{eq:angle3d}
\end{equation}

Durch die Normierung mit den Beträgen der beiden Vektoren geht nur deren Richtung und nicht deren Länge in das Ergebnis ein. Gelenkwinkel sind deshalb unabhängig von einer einheitlichen Skalierung des Körpers. Tabelle~\ref{tab:angles} fasst die verwendeten Winkelgrößen zusammen.

\begin{table}[h]
\centering
\small
\begin{tabular}{p{0.28\textwidth} p{0.38\textwidth} p{0.13\textwidth}}
\hline
\textbf{Größe} & \textbf{Definition} & \textbf{Dimension} \\ \hline
Kniewinkel & Hüfte -- \underline{Knie} -- Sprunggelenk & 3D \\
Hüftbeugung & Schulter -- \underline{Hüfte} -- Knie & 3D \\
Oberschenkelneigung & Vektor Hüfte$\rightarrow$Knie gegen Vertikale $(0,1,0)$ & 3D \\
Körperrelative Hüftausrichtung & Beckenrichtung gegen Richtung Hüfte$\rightarrow$Sprunggelenk des Kickbeins & 2D in der $x$-$z$ Ebene \\ \hline
\end{tabular}
\caption{Verwendete Winkelgrößen und ihre Definition}
\label{tab:angles}
\end{table}

Die körperrelative Hüftausrichtung wird im Gegensatz zu den Gelenkwinkeln in der Draufsicht auf die $x$-$z$ Ebene bestimmt. Zunächst wird aus der Verbindungslinie beider Hüften eine dazu senkrechte Beckenrichtung abgeleitet. Diese wird mit dem Vektor von der Hüfte zum Sprunggelenk des Kickbeins verglichen.

Wichtig ist die Interpretation dieses Winkels. Er beschreibt nicht die globale Drehung des Körpers im Raum, sondern die Lage des Kickbeins relativ zum Becken. Drehen sich Becken und Kickbein gemeinsam um die Körperlängsachse, kann der körperrelative Winkel unverändert bleiben. Die drei Techniken werden deshalb über unterschiedliche relative Zielausrichtungen beschrieben (Tabelle~\ref{tab:hip_alignment_technique}).

\begin{table}[h]
\centering
\begin{tabular}{lll}
\hline
\textbf{Kickbein relativ zum Becken} & \textbf{Wert} & \textbf{Technik} \\ \hline
annähernd in Beckenrichtung & $\sim 0°$ & Front Kick \\
seitlich & $\sim 90°$ & Roundhouse Kick \\
nach hinten & $\sim 180°$ & Spinning Back Kick \\ \hline
\end{tabular}
\caption{Interpretation der körperrelativen Hüftausrichtung je Technik}
\label{tab:hip_alignment_technique}
\end{table}

Alle drei Techniken werten diese Größe im Bereich von 0 bis 180 Grad aus. Gegen sprunghafte Vertauschungen der Hüftlandmarken wird der in Abschnitt~4.3.4 beschriebene Plausibilitätsfilter eingesetzt. Eine Änderung von mehr als 60 Grad zwischen zwei aufeinanderfolgenden Bildern wird verworfen. Dieser Filter reduziert offensichtliche Sprünge, ersetzt jedoch keine unabhängige Prüfung der zugrunde liegenden Landmark-Qualität.

\subsection{Kriterien und Schwellenwerte je Technik}
Tabelle~\ref{tab:crit_front} bis Tabelle~\ref{tab:crit_dwit} listen die Bewertungskriterien der drei Techniken mit den zugehörigen \texttt{fail\_at} und \texttt{ideal\_at} Werten aus Gleichung~\ref{eq:score}. Die beiden Werte definieren die Endpunkte der linearen Abbildung auf 0 bis 100 Prozent. Sie sind nicht mit der Bestehensgrenze gleichzusetzen, die in der Benutzeroberfläche bei 50 Prozent liegt.

\begin{table}[htbp]
\centering
\small
\begin{tabular}{c p{0.18\textwidth} p{0.34\textwidth} c c}
\hline
\# & \textbf{Kriterium} & \textbf{Messgröße} & \textbf{fail\_at} & \textbf{ideal\_at} \\ \hline
1 & Beinstreckung & maximaler Kniewinkel, Kick-Phase & 131° & 170° \\
2 & Chamber-Winkel & minimaler Kniewinkel, Chamber-Phase & 100° & 71° \\
3 & Hüftbeugung & maximale Oberschenkelneigung, Chamber-Phase & 50° & 90° \\
4 & Hüftausrichtung & körperrelative Hüftausrichtung im Treffmoment & 90° & 30° \\
5 & Knieabsinken & Absinken des Knies während der Kick-Phase & 0,30\,m & 0,15\,m \\
6 & Standbein & Median des Standbein-Kniewinkels, Kick-Phase & 180° & 165° \\
7 & Rechamber & Rückzug vom Streckungsmaximum & 10° & 60° \\ \hline
\end{tabular}
\caption{Bewertungskriterien des \textit{Front Kicks}}
\label{tab:crit_front}
\end{table}

Für den \textit{Front Kick} lassen sich zwei Grenzwerte an publizierter Kinematik orientieren. Lee et al. \cite{Lee2005} berichten bei vier erfahrenen Athleten einen Kniewinkel von 71,0~$\pm$~1,33 Grad im Moment der stärksten Beugung und 131,2~$\pm$~2,25 Grad im Treffmoment.

Für den Chamber-Winkel wird der Wert von 71 Grad als \texttt{ideal\_at} übernommen, da er der in dieser Stichprobe gemessenen starken Beugung entspricht. Als \texttt{fail\_at} wird dagegen die Schwelle von 100 Grad verwendet, die zugleich den Eintritt in die Chamber-Phase begrenzt. Diese Grenze ist keine publizierte Qualitätsgrenze, sondern eine prototypische Setzung. Sie lässt auch weniger stark angewinkelte Ausführungen noch in die Bewertung eingehen.

Für die Beinstreckung dient der von Lee et al. berichtete Treffmoment-Wert von 131,2 Grad als Orientierung für die untere Bewertungsgrenze. Die Übertragung ist nur eingeschränkt möglich, da Lee et al. einen tatsächlichen Treffer untersuchten, während die vorliegende Anwendung Kicks ohne physisches Ziel auswertet. Der Wert wird daher nicht als allgemeingültige Grenze korrekter Technik verstanden. \texttt{ideal\_at} liegt bei 170 Grad, damit eine weitergehende Streckung bei Kicks in die Luft höhere Punktwerte erreichen kann. Beide Grenzwerte müssen in einer späteren fachlichen Validierung überprüft werden.

Die übrigen Grenzwerte des \textit{Front Kicks} wurden aus der Trainingspraxis und der prototypischen Entwicklung abgeleitet. Dies gilt auch für das Standbein. Die aktuelle Parametrisierung vergibt bei einem Kniewinkel von 165 Grad die volle Punktzahl und bei 180 Grad null Punkte. Sie bevorzugt damit eine leichte Beugung gegenüber einer vollständigen Streckung. Da hierfür keine empirische Kalibrierung vorliegt, muss gerade diese Bewertungsrichtung in einer späteren Validierung überprüft werden.

\begin{table}[htbp]
\centering
\small
\begin{tabular}{c p{0.18\textwidth} p{0.34\textwidth} c c}
\hline
\# & \textbf{Kriterium} & \textbf{Messgröße} & \textbf{fail\_at} & \textbf{ideal\_at} \\ \hline
1 & Beinstreckung & maximaler Kniewinkel, Kick-Phase & 131° & 170° \\
2 & Hüftausrichtung & Abweichung der körperrelativen Hüftausrichtung im Treffmoment von 90° & 60° & 20° \\
3 & Chamber-Winkel & minimaler Kniewinkel, Chamber-Phase & 100° & 71° \\
4 & Hüftbeugung & maximale Oberschenkelneigung, Chamber-Phase & 50° & 90° \\
5 & Knieabsinken & Absinken des Knies während der Kick-Phase & 0,30\,m & 0,15\,m \\
6 & Standbein & Median des Standbein-Kniewinkels, Kick-Phase & 180° & 165° \\
7 & Rechamber & Rückzug vom Streckungsmaximum & 10° & 60° \\ \hline
\end{tabular}
\caption{Bewertungskriterien des \textit{Roundhouse Kicks}}
\label{tab:crit_round}
\end{table}

Ein wesentliches technisches Merkmal des \textit{Roundhouse Kicks} ist die Beckenrotation. Die vorliegende Implementierung misst die globale Rotation des Beckens jedoch nicht direkt. Bewertet wird stattdessen die körperrelative Ausrichtung des Kickbeins zum Becken im Treffmoment (Abschnitt~4.7). Dieses Maß dient als Näherung für die technikspezifische Ausrichtung im Moment der stärksten Beinstreckung.

Beim \textit{Front Kick} wird ein kleiner Winkel angestrebt. Beim \textit{Roundhouse Kick} liegt der Zielbereich dagegen um 90 Grad. Weil sowohl eine geringere als auch eine stärkere Abweichung ungünstig sein kann, wird für die Bewertung nicht der Rohwert verwendet, sondern seine absolute Abweichung von 90 Grad:

\begin{equation}
\theta_{\text{Fehler}} = \left|\text{hip\_alignment\_at\_impact} - 90^{\circ}\right|
\label{eq:rotationerror}
\end{equation}

Diese Abweichung geht mit \texttt{fail\_at}=60 Grad und \texttt{ideal\_at}=20 Grad in Gleichung~\ref{eq:score} ein. Angezeigt wird zusätzlich der zugrunde liegende Rohwert. Ein Wert von 70 Grad und einer von 110 Grad liegen jeweils 20 Grad vom Zielwert entfernt und erhalten deshalb dieselbe Punktzahl. Die Richtung der Abweichung kann im Rückmeldungstext unterschieden werden, nicht jedoch in der Punktzahl selbst.

\begin{table}[htbp]
\centering
\small
\begin{tabular}{c p{0.18\textwidth} p{0.34\textwidth} c c}
\hline
\# & \textbf{Kriterium} & \textbf{Messgröße} & \textbf{fail\_at} & \textbf{ideal\_at} \\ \hline
1 & Beinstreckung & maximaler Kniewinkel, Kick-Phase & 131° & 170° \\
2 & Chamber-Winkel & minimaler Kniewinkel, Chamber-Phase & 100° & 71° \\
3 & Beinausrichtung & körperrelative Hüftausrichtung im Treffmoment & 110° & 170° \\
4 & Standbein & Median des Standbein-Kniewinkels, Kick-Phase & 180° & 165° \\
5 & Rechamber & Rückzug vom Streckungsmaximum & 10° & 60° \\ \hline
\end{tabular}
\caption{Bewertungskriterien des \textit{Spinning Back Kicks}}
\label{tab:crit_dwit}
\end{table}

Beim \textit{Spinning Back Kick} bleibt der Kniewinkel für die Bewertung relevant, obwohl er für den Übergang in die Kick-Phase nicht als alleiniges Erkennungsmerkmal verwendet wird. Kim et al. \cite{Kim2011} untersuchten bei zwölf Taekwondo-Athleten mit Schwarzgurt vier Kicktechniken. Für den \textit{Spinning Back Kick} berichten sie den größten Bewegungsumfang des Knies und beschreiben die Technik als \textit{push-like movement}, bei der Hüfte und Knie zur Erzeugung der Kickgeschwindigkeit gemeinsam strecken. Diese Bewegungscharakteristik begründet, weshalb sowohl eine deutliche Beugung als auch eine ausgeprägte Streckung in der Bewertung berücksichtigt werden.

Für die Beinstreckung wird derselbe Bereich wie bei den beiden anderen Techniken verwendet. Anders als beim \textit{Front Kick} existiert für den \textit{Spinning Back Kick} jedoch kein direkt vergleichbarer Treffmoment-Wert aus der herangezogenen Literatur. Die Grenzwerte von 131 und 170 Grad sind hier deshalb als prototypische, expertenbasierte Setzungen zu verstehen und nicht als empirisch bestätigte Referenzwerte.

Das Kriterium \textit{Beinausrichtung} bewertet die relative Lage des Kickbeins zum Becken und darf nicht mit der globalen Körperdrehung gleichgesetzt werden. Ein Wert in Richtung 180 Grad bedeutet, dass das Kickbein im Verhältnis zur abgeleiteten Beckenrichtung nach hinten ausgerichtet ist. Die volle Punktzahl wird bereits bei 170 Grad vergeben, da 180 Grad die mathematische Obergrenze des verwendeten Winkelmaßes darstellt und Messfehler in Grenznähe die volle Punktzahl sonst unnötig empfindlich machen würden. Mit \texttt{fail\_at}=110 Grad wird eine Ausrichtung unterhalb des für diese Technik vorgesehenen hinteren Bereichs mit 0 Prozent bewertet.

Zusätzlich wird beim \textit{Spinning Back Kick} eine aufsummierte Nachdrehung geführt. Sie entscheidet lediglich darüber, welcher Rückmeldungstext zur Eindrehung angezeigt wird. Wegen der in Kapitel~6 beschriebenen Einschränkungen der Rotationszählung geht diese Hilfsgröße nicht in die Punktzahl ein.

Ein Vergleich mit Wang et al. (Abschnitt~3.1) zeigt Überschneidungen bei den Bereichen Standbein, Kicktechnik, Chamber und Rückzug. Die konkrete Berechnung und Parametrisierung unterscheidet sich jedoch. Landeposition und Handhaltung werden in der vorliegenden Arbeit nicht bewertet, da die dafür relevanten Landmarken an Füßen und Handgelenken nicht ausgewertet werden.

\subsection{Abbruchbedingungen bei der Bewertung}
Nicht jedes Video enthält einen vollständig ausgeführten Kick. Alle drei Techniken brechen die Bewertung vorzeitig ab, wenn im Video nie eine Kick-Phase erreicht wird. Ausschlaggebend ist dafür das Merkmal \texttt{kick\_detected}, das beim ersten Wechsel in die Kick-Phase gesetzt wird. Bleibt es auf \texttt{False}, berechnet \texttt{evaluate()} nur das erste Kriterium, also die Beinstreckung und kehrt danach sofort zurück. Die übrigen Kriterien entfallen.

Für dieses eine Kriterium ergibt sich dann zwangsläufig ein Wert von 0 Prozent. Der zugehörige Messwert \texttt{max\_knee\_in\_kick} wird nämlich nur während der Kick-Phase fortgeschrieben und behält ohne sie seinen Startwert 0,0. Da der Gesamtwert der Mittelwert über alle zurückgegebenen Kriterien ist (Abschnitt~4.10), fällt auch dieser auf 0.
Das Ergebnis ist in dieser Version nicht immer angemessen. Verfehlt eine Ausführung die Erkennungsschwelle nur knapp, wird sie genauso mit 0 Prozent bewertet wie ein Video, in dem überhaupt nicht getreten wurde. Aus diesem Grund wurde die Schwelle für den Übergang in die Kick-Phase bewusst niedrig angesetzt (Abschnitt~4.4.4). Beim \textit{Front Kick} und beim \textit{Roundhouse Kick} liegt sie mit 131 Grad genau auf dem \texttt{fail\_at}-Wert der Beinstreckung, sodass eine gerade noch erkannte Ausführung bei 0 Prozent beginnt und mit zunehmender Streckung mehr Punkte erhält. 
Beim \textit{Spinning Back Kick} wird die Kick-Phase dagegen anhand der Fußhöhe und des Knöchelabstands erkannt. Der Kniewinkel ist für die Erkennung hier nicht ausschlaggebend, sodass eine Ausführung mit einem Kniewinkel unter 131 Grad trotzdem erkannt und bewertet werden kann. Das grundsätzliche Problem bleibt jedoch bestehen: Wird die Erkennungsschwelle nicht erreicht, wird keine abgestufte Bewertung vorgenommen, sondern die Ausführung mit 0 Prozent bewertet. Die Schwelle bestimmt damit weiterhin, ob eine Bewertung stattfindet, und nicht nur, wie gut die Ausführung ist. Dieser Zielkonflikt wird in Kapitel~6 erneut aufgegriffen.
\subsection{Gesamtergebnis}
Der angezeigte Gesamtwert ist der ungewichtete arithmetische Mittelwert über alle $n$ zurückgegebenen Kriterienwerte:
\begin{equation}
\text{overall} = \frac{1}{n} \sum_{i=1}^{n} \text{score}_i
\label{eq:overall}
\end{equation}
Aus dem Gesamtwert wird eine Sternebewertung zwischen 1 und 5 abgeleitet:
\begin{equation}
\text{stars} = \min\!\left(5,\ \max\!\left(1,\ \operatorname{round}\!\left(\frac{\text{overall}}{20}\right)\right)\right)
\label{eq:stars}
\end{equation}
Die verbale Einordnung erfolgt in vier Stufen: \glqq{}Sehr gut!\grqq{} ab 90 Prozent, \glqq{}Gut!\grqq{} ab 75 Prozent, \glqq{}Solide\grqq{} ab 50 Prozent, andernfalls \glqq{}Übungsbedarf\grqq{}. Die Tipps-Liste enthält die Rückmeldungstexte aller Kriterien, deren Wert unter 50 Prozent liegt.

Eine Gewichtung der Kriterien findet nicht statt. Alle Kriterien gehen gleich stark in Gleichung~\ref{eq:overall} ein. Dadurch kann eine falsch gewählte Technik das Gesamtergebnis nur begrenzt beeinflussen. Wird beispielsweise ein sauber ausgeführter \textit{Roundhouse Kick} als \textit{Spinning Back Kick} bewertet, beträgt die Hüftausrichtung im Treffmoment etwa 90° statt der angestrebten 170°. Damit liegt sie unterhalb des \texttt{fail\_at}-Wertes von 110° und erhält für die Körperdrehung 0 Prozent. Die übrigen vier Kriterien bewerten Chamber, Streckung, Standbein und Rückzug und können weiterhin hohe Werte erreichen. Dadurch sind theoretisch bis zu 80 Prozent möglich. Ein Testvideo erreicht in diesem Fall rund 78 Prozent. Auch diese Einschränkung wird in Kapitel~6 aufgegriffen.
\subsection{Geschwindigkeitsmessung}
Die relative Fußgeschwindigkeit wird aus der Positionsänderung des Kickknöchels zwischen zwei aufeinanderfolgenden Bildern berechnet:
\begin{equation}
v_{\text{Fuß}} = d \cdot f \cdot 3{,}6 \quad [\text{km/h}]
\label{eq:speed}
\end{equation}
Die Variable $d$ ist die euklidische Distanz zwischen zwei aufeinanderfolgenden Knöchelpositionen in Metern und $f$ die Bildrate in Bildern pro Sekunde. Der Faktor $3{,}6$ rechnet von Metern pro Sekunde in Kilometer pro Stunde um. Der so berechnete Wert wird über drei Bilder medianglättet und als Maximum über Chamber- und Kick-Phase geführt.

Die Weltkoordinaten von MediaPipe sind hüftzentriert (Abschnitt~2.3.3). Gleichung~\ref{eq:speed} misst deshalb die Fußbewegung relativ zur Hüfte und nicht absolut im Raum. Verlagert sich der ganze Körper nach vorn, etwa durch ein Vorspringen beim Kick, fehlt dieser Anteil im Ergebnis. Der Wert beschreibt damit keine vollständige Geschwindigkeit des Fußes im Raum, sondern die Bewegung relativ zur Hüftmitte. Er dient allein der Anzeige und fließt in kein Bewertungskriterium ein.
\subsection{Nutzerschnittstelle und Rückmeldung}
Die Technik wird vor der Analyse vom Nutzer ausgewählt, eine automatische Erkennung findet nicht statt. Zusätzlich lässt sich die Modellvariante zwischen \textit{Lite, Full} und \textit{Heavy} umschalten. Für eine Endanwendung wäre diese Wahl ungewöhnlich, sie ist als Eigenschaft des Prototyps zu verstehen und erlaubt es, den Einfluss der Modellgröße systematisch zu untersuchen (Abschnitt~4.2.1).

\begin{figure}[h]
    \centering
    \includegraphics[width=0.6\textwidth]{Images/app\_technikwahl.png}
    \caption{Technik- und Modellauswahl in der Weboberfläche vor dem Videoupload.}
    \label{fig:app_technikwahl}
\end{figure}

Nach Abschluss der Analyse liefert die Anwendung zwei Ausgaben. Die erste ist ein annotiertes Video. Es zeigt das Skelett aus den 33 \textit{Landmarks} mit grün eingefärbten Verbindungslinien, wobei Punkt- und Liniendicke mit der Videoauflösung skalieren.

Darüber liegt ein Debug-Overlay mit den aktuellen Messwerten: Phase, Kniewinkel, Hüftbeugung, Kickseite, Oberschenkelneigung, Startbild, Kickdauer, Sichtbarkeit und Hüftausrichtung. Beim \textit{Spinning Back Kick} kommen die Ausrichtung im Treffmoment, die Nachdrehung, die Gesamtdrehung, der Fußabstand und die Fußhöhe hinzu.

\begin{figure}[h]
    \centering
    \includegraphics[width=0.6\textwidth]{Images/app\_video\_overlay.png}
    \caption{Annotiertes Ausgabevideo mit Skelett- und Debug-Overlay, hier im Moment der Kickstreckung.}
    \label{fig:app_video_overlay}
\end{figure}

Die zweite Ausgabe ist die Ergebniskarte. Sie zeigt den Gesamtwert als Zahl von 0 bis 100, die relative Fußgeschwindigkeit in km/h, die verbale Einordnung und die Sternebewertung. Darunter steht für jedes Kriterium eine Zeile mit Bezeichnung, Balken und Prozentwert. Die Balken sind ab 90 Prozent grün, ab 75 Prozent gelb und darunter rot. Die Farbgrenzen stimmen damit nicht vollständig mit der Bestehensgrenze von 50 Prozent überein. Ein Kriterium zwischen 50 und 74 Prozent wird bereits als bestanden behandelt und erzeugt keinen Verbesserungstipp, erscheint jedoch weiterhin rot. Diese Inkonsistenz der Rückmeldelogik ist eine gestalterische Grenze des Prototyps und sollte in einer Weiterentwicklung vereinheitlicht werden. Am Ende listet die Anwendung die Rückmeldungen aller Kriterien unter 50 Prozent auf (Abschnitt~4.10).

\begin{figure}[h]
    \centering
    \includegraphics[width=0.45\textwidth]{Images/ergebniskarte\_gelb.png}\hfill
    \includegraphics[width=0.45\textwidth]{Images/ergebniskarte\_gruen.png}
    \caption{Ergebniskarten mit Gesamtwert, Sternebewertung, Kriterienliste und Tipps.}
    \label{fig:app_ergebniskarte}
\end{figure}
Vor der ersten Aufnahme erscheint automatisch ein Hinweistext zu den Aufnahmebedingungen (Abschnitt~4.14), der über eine Schaltfläche erneut aufrufbar ist.

\begin{figure}[h]
    \centering
    \includegraphics[width=0.6\textwidth]{Images/app\_filmingtips.png}
    \caption{Hinweisdialog zu den Aufnahmebedingungen vor der Analyse.}
    \label{fig:app_filmingtips}
\end{figure}
\subsection{Weitere Funktionen}
Analysierte Ausführungen lassen sich speichern. Das Video wird dabei als MP4 abgelegt, die Bewertung als JSON. Der Verlauf ist pro Browserprofil getrennt. Dafür erzeugt das Frontend beim ersten Aufruf eine zufällige Kennung im \texttt{localStorage} des Browsers, schreibt sie beim Speichern mit und verwendet sie beim Abruf als Filter.

Diese Kennung dient allein der Zuordnung gespeicherter Analysen zu einem Browserprofil. Wer den Browser wechselt oder die Browserdaten löscht, verliert den Zugang zum bisherigen Verlauf. Für einen Prototyp, der die eigenen Ausführungen verarbeitet, ist das vertretbar, für einen Einsatz mit mehreren Personen wäre eine Authentifizierung erforderlich.

\begin{figure}[h]
    \centering
    \includegraphics[width=0.6\textwidth]{Images/app\_verlauf.png}
    \caption{Verlaufsseite mit gespeicherten Analysen eines Geräts.}
    \label{fig:app_verlauf}
\end{figure}

Ergänzend zur Bewegungsanalyse enthält die Anwendung eine nach Gurtgrad gegliederte Vokabelliste sowie eine Vokabelabfrage im Multiple-Choice-Format. Abgedeckt sind die koreanischen Fachbegriffe der Gürtelprüfungen, die in Abschnitt~1.3 genannt werden.

\begin{figure}[h]
    \centering
    \includegraphics[width=0.45\textwidth]{Images/app\_vokabeln.png}\hfill
    \includegraphics[width=0.45\textwidth]{Images/app\_vokabeltest.png}
    \caption{Gürtelabhängige Vokabelliste und Multiple-Choice-Abfrage.}
    \label{fig:app_vokabeln}
\end{figure}
\subsection{Aufnahmebedingungen}
Die Anforderungen an die Aufnahme werden ausschließlich als Hinweistext vermittelt (Abbildung~\ref{fig:app_filmingtips}). Gefordert werden eine Aufnahme von der Seite oder schräg im 45-Grad-Winkel, nur eine Person im Bild, ein auch im höchsten Punkt des Kicks vollständig sichtbarer Körper, ein Abstand von etwa drei bis vier Metern, eine möglichst ruhige Kamera, nur ein Kick pro Video, gute Beleuchtung, ein ruhiger Hintergrund und eine Hochkant-Aufnahme. 
Technisch erzwungen wird davon nichts. Perspektive, Abstand, Beleuchtung und Hintergrund werden nicht geprüft. Befinden sich mehrere Personen im Bild, wertet das System die von MediaPipe je Bild gelieferte Pose aus, die dabei zwischen den Bildern wechseln kann. Lediglich die Bildrate wird aus den Metadaten gelesen, mit einem Standardwert bei 30 Bildern pro Sekunde (Abschnitt~4.3.1).

Lässt sich die Datei nicht öffnen, wird der Auftragsstatus auf \textit{error} gesetzt und die Meldung \glqq{}Konnte Video nicht öffnen.\grqq{} ausgegeben. Wird keine Person erkannt, werden schlicht keine Messwerte erhoben.
\section{Evaluation \& Bewertung}
\newcommand{\evph}[1]{[\textit{#1}]}
\newcommand{\Nteiln}{19}
\newcommand{\Nfreitext}{16}
\newcommand{\AlterMittel}{29{,}9}
\newcommand{\AlterSD}{8{,}8}
\newcommand{\AlterMedian}{27}
\newcommand{\AlterMin}{21}
\newcommand{\AlterMax}{49}
\newcommand{\SUSm}{88{,}0}
\newcommand{\SUSsd}{9{,}5}
\newcommand{\SUSmd}{90{,}0}
\newcommand{\SUSmin}{67{,}5}
\newcommand{\SUSmax}{100}

\subsection{Zielsetzung und Aufbau der Evaluation}
\label{sec:eval_ziel}
Die Nutzerstudie untersucht die zweite Forschungsfrage aus Abschnitt~1.3. Im Mittelpunkt stehen damit die wahrgenommene Gebrauchstauglichkeit, die Verständlichkeit der Rückmeldung, das Vertrauen in die Bewertung, der wahrgenommene Nutzen und die Nutzungsabsicht. Ergänzend werden Erfahrungen mit der technischen Personenerkennung, der Vokabelfunktion und der Gestaltung erhoben.

Die erste Forschungsfrage betrifft die technische Umsetzung und die Grenzen der regelbasierten Analyse. Kapitel~4 zeigt, wie die ausgewählten Messgrößen und Kriterien umgesetzt wurden. Eine unabhängige messtechnische Validierung der berechneten Winkel und Positionen sowie ein systematischer Vergleich der Bewertung mit fachlichen Trainerurteilen wurden jedoch nicht durchgeführt. Die Nutzerstudie kann deshalb keine Aussage darüber treffen, ob die Messwerte objektiv richtig sind oder ob die Punktwerte die tatsächliche technische Qualität eines Kicks korrekt abbilden. Sie zeigt ausschließlich, wie die Zielgruppe den Prototyp und seine Rückmeldungen wahrnimmt.

Diese Trennung ist für die Interpretation der Ergebnisse zentral. Eine hohe Akzeptanz oder ein hohes Vertrauen kann die fachliche Validität der Bewertung nicht ersetzen. Umgekehrt kann eine technisch korrekte Berechnung für den Trainingsalltag wenig hilfreich sein, wenn ihre Rückmeldung nicht verstanden wird. Die vorliegende Evaluation untersucht daher gezielt die nutzerbezogene Seite des Systems. Möglichkeiten für eine spätere technische und fachliche Validierung werden in Abschnitt~\ref{sec:ausblick} beschrieben.

\subsection{Hypothesen}
\label{sec:eval_hyp}
Für die Nutzerstudie werden sechs Hypothesen formuliert. H1 bezieht sich auf die allgemeine Gebrauchstauglichkeit. H2 bis H5 betreffen die im Fragebogen erhobenen Bereiche Aufnahme, Verständlichkeit, Vertrauen und wahrgenommenen Nutzen. H6 erfasst die Nutzungsabsicht.

\begin{table}[h]
\centering
\small
\begin{tabular}{p{0.05\textwidth} p{0.50\textwidth} p{0.27\textwidth}}
\hline
& \textbf{Aussage} & \textbf{Datenquelle} \\ \hline
H1 & Der mittlere SUS-Wert liegt über dem Vergleichswert von 68 & SUS-Block \\
H2 & Die Aufnahme ist einfach und verständlich durchzuführen & C1 Aufnahme \\
H3 & Die Bewertung ist für die Nutzenden verständlich und nachvollziehbar & C2 Verständlichkeit \\
H4 & Die Nutzenden vertrauen der Bewertung & C3 Vertrauen \\
H5 & Die Anwendung wird als nützlich für das eigene Training wahrgenommen & C4 Nützlichkeit \\
H6 & Es besteht die Absicht, die Anwendung zur Prüfungsvorbereitung einzusetzen oder weiterzuempfehlen & C5 Nutzungsabsicht \\ \hline
\end{tabular}
\caption{Hypothesen der Nutzerstudie und ihre Datengrundlage}
\label{tab:eval_hyp}
\end{table}

Die Hypothesen werden deskriptiv eingeordnet. H1 wird mit dem in der SUS-Literatur häufig verwendeten Vergleichswert von 68 verglichen \cite{SauroLewis2016}. H2 bis H5 werden anhand der Mittelwerte, Mediane und Zustimmungsanteile der zugehörigen Aussagen beurteilt. Für H6 liegen keine fünfstufigen Ratings vor, sondern eine Mehrfachauswahl. Deshalb werden dort ausschließlich die Auswahlhäufigkeiten berichtet. Die Hypothesen dienen der strukturierten Beschreibung der Nutzerstudie und nicht der Prüfung kausaler Zusammenhänge.

\subsection{Auswertungsverfahren}
\label{sec:eval_auswertung}
Die Auswertung erfolgt deskriptiv. Für jede fünfstufige Aussage werden Mittelwert, Median und Zustimmungsanteil berichtet. Als Zustimmung gelten die beiden oberen Skalenstufen. Die Hypothesen H2 bis H5 werden anhand der Richtung und Höhe dieser Verteilungen eingeordnet. H1 wird anhand des SUS-Gesamtwerts beurteilt, H6 anhand der Auswahlhäufigkeiten.

Auf inferenzstatistische Signifikanztests wird verzichtet. Die Stichprobe ist eine kleine Gelegenheitsstichprobe ohne Zufallsauswahl und ohne Kontrollgruppe. Ziel der Untersuchung ist die explorative Beurteilung des entwickelten Prototyps durch die konkrete Zielgruppe. Die Ergebnisse werden deshalb nicht als repräsentative Schätzung für alle Taekwondo-Trainierenden interpretiert.

Negativ formulierte Aussagen werden nicht umkodiert, sondern mit ihren Rohwerten berichtet und in den Tabellen mit einem Minuszeichen gekennzeichnet. Bei ihnen bedeutet ein niedriger Wert ein positives Ergebnis. Auf Skalenmittelwerte über mehrere selbst formulierte Aussagen hinweg wird verzichtet, da die einzelnen Aussagen unterschiedliche Aspekte des jeweiligen Bereichs erfassen.

Eine Ausnahme bildet die \textit{System Usability Scale}. Sie ist ein etabliertes Instrument mit einer festgelegten Berechnungsvorschrift \cite{Brooke1996}. Die zehn Items werden zu einem Gesamtwert zwischen 0 und 100 verrechnet. Dieser Wert ist nicht als Prozentwert zu interpretieren. Für die Einordnung wird ein Wert von 68 als häufig verwendeter Vergleichswert herangezogen \cite{SauroLewis2016}.

\subsection{Methodisches Vorgehen der Nutzerstudie}
\label{sec:eval_methode}
Die Erhebung fand im laufenden Taekwondo-Training statt. Die Teilnehmenden wurden in Zweiergruppen eingeteilt und probierten die Anwendung gemeinsam aus. Die Zweiergruppen ergaben sich aus dem Aufnahmeszenario. Während eine Person den Kick ausführte, konnte die zweite Person filmen. Anschließend konnten die Rollen getauscht werden.

Vor der Nutzung erhielten die Teilnehmenden eine kurze mündliche Einweisung in Zweck und Grundidee der Anwendung. Konkrete Bedienschritte, die Bedeutung einzelner Bewertungskriterien und zusätzliche Aufnahmehinweise wurden nicht erklärt. Diese Informationen sollten soweit möglich aus der Anwendung selbst hervorgehen.

Die Gruppen konnten die Funktionen anschließend frei ausprobieren. Es gab weder eine fest vorgegebene Aufgabenfolge noch eine einheitliche Nutzungsdauer. Diese offene Durchführung entspricht einer explorativen Erprobung, reduziert jedoch die Vergleichbarkeit zwischen den Teilnehmenden. Zudem konnten sich die Personen innerhalb einer Zweiergruppe während der Nutzung gegenseitig beeinflussen. Der Fragebogen wurde anschließend von jeder Person einzeln ausgefüllt.

Der Fragebogen war anonym und die Teilnahme freiwillig. Er gliedert sich in acht Blöcke: Angaben zur Person und Trainingserfahrung, Nutzungskontext, Kickanalyse mit Aussagen zu Aufnahme, Verständlichkeit, Vertrauen und Nützlichkeit, technische Erkennung, Vokabelfunktion, Design, die System Usability Scale sowie offene Fragen. Die Aussagenblöcke wurden auf einer fünfstufigen Skala von \glqq{}Stimme überhaupt nicht zu\grqq{} bis \glqq{}Stimme voll zu\grqq{} beantwortet.

Ein Teil der Aussagen ist negativ formuliert, beispielsweise \glqq{}Ich hatte den Eindruck, dass die App meinen Kick teilweise falsch bewertet\grqq{}. In den folgenden Tabellen sind diese Aussagen mit einem Minuszeichen gekennzeichnet. Bei ihnen bedeutet ein niedriger Wert ein positives Ergebnis.

\subsection{Stichprobe}
\label{sec:eval_stichprobe}
An der Befragung nahmen \Nteiln{} Personen teil. Das Durchschnittsalter beträgt \AlterMittel{} Jahre bei einer Standardabweichung von \AlterSD{} Jahren. Der Median liegt bei \AlterMedian{} Jahren, die Altersspanne reicht von \AlterMin{} bis \AlterMax{} Jahren.

Für die Einordnung der Ergebnisse sind zwei Merkmale besonders relevant. Erstens liegt der Schwerpunkt der Stichprobe bei den unteren Gurtgraden. Fünfzehn der \Nteiln{} Personen tragen Weiß, Gelb, Gelb Grün oder Grün. Die Beurteilung stammt damit überwiegend von Trainierenden, die selbst noch beim Erlernen der Techniken sind, und nicht von erfahrenen Prüferinnen und Prüfern. Zweitens hatten zwölf Personen zuvor noch nie Videoaufnahmen zur Kontrolle der eigenen Technik genutzt. Für diese Personen war neben der Anwendung auch das Prinzip der Videoanalyse neu.

Alle \Nteiln{} Personen gaben an, sämtliche Funktionen ausprobiert zu haben. Den \textit{Front Kick} testeten alle, den \textit{Roundhouse Kick} 17 und den \textit{Spinning Back Kick} 14 Personen. Die rotationsstärkste und technisch anspruchsvollste der drei betrachteten Techniken ist damit zugleich am schwächsten vertreten. Tabelle~\ref{tab:eval_stichprobe} fasst die Stichprobe zusammen.

\begin{table}[htbp]
\centering
\small
\begin{tabular}{p{0.34\textwidth} p{0.50\textwidth}}
\hline
\textbf{Merkmal} & \textbf{Verteilung} \\ \hline
Alter & \AlterMittel{} Jahre ($\pm$~\AlterSD{}), Spanne \AlterMin{} bis \AlterMax{} \\
Geschlecht & 10 männlich, 9 weiblich \\
Gürtelgrad & 1 Weiß, 5 Gelb, 5 Gelb Grün, 4 Grün, 1 Blau Rot, 1 Rot, 1 Rot Schwarz, 1 Schwarz \\
Trainingsdauer & 2 unter 1 Jahr, 10 ein bis drei Jahre, 5 drei bis fünf Jahre, 2 über 5 Jahre \\
Trainingshäufigkeit & 7 einmal, 10 zweimal, 2 dreimal pro Woche \\
Vorerfahrung mit Videoanalyse & 7 gelegentlich, 12 keine \\
Getestete Techniken & 19 \textit{Front Kick}, 17 \textit{Roundhouse Kick}, 14 \textit{Spinning Back Kick} \\ \hline
\end{tabular}
\caption{Zusammensetzung der Stichprobe ($n=\Nteiln$)}
\label{tab:eval_stichprobe}
\end{table}

\subsection{Ergebnisse zur Kickanalyse}
\label{sec:eval_kick}
Die folgenden Tabellen führen die Aussagen der vier Bereiche Aufnahme, Verständlichkeit, Vertrauen und Nützlichkeit mit Mittelwert, Median und Zustimmungsanteil auf.

\begin{table}[htbp]
\centering
\small
\begin{tabular}{p{0.46\textwidth} c c c}
\hline
\textbf{Aussage} & \textbf{M} & \textbf{Md} & \textbf{Zust.} \\ \hline
Das Aufnehmen der Videos war unkompliziert & 4,74 & 5 & 100\,\% \\
Mir war klar, wie ich die Kamera positionieren musste & 4,47 & 5 & 89\,\% \\
Ich war mir unsicher, ob ich richtig im Bild stehe ($-$) & 2,47 & 2 & 21\,\% \\ \hline
\end{tabular}
\caption{Aussagen zur Aufnahmesituation ($n=\Nteiln$)}
\label{tab:eval_aufnahme}
\end{table}

Die Aufnahmesituation wurde überwiegend positiv beurteilt. Alle Teilnehmenden stimmten der Aussage zu, dass das Aufnehmen unkompliziert war. 89 Prozent gaben an, dass ihnen die Positionierung der Kamera klar war. Vier Personen stimmten der negativ formulierten Aussage zu, sie seien unsicher gewesen, ob sie richtig im Bild stehen. H2 wird damit durch die Antworten gestützt.

Unabhängig davon berichteten nur sechs Personen, dass kein Video erneut aufgenommen werden musste. Neun Personen nahmen einmal neu auf und vier Personen zweimal. Insgesamt wiederholten damit 13 der \Nteiln{} Teilnehmenden mindestens eine Aufnahme. Aus dem Fragebogen lässt sich jedoch nicht eindeutig ableiten, weshalb die Wiederholung erforderlich war. Gründe können sowohl in der Aufnahme, in der Personenerkennung als auch in der ausgeführten Technik oder im Wunsch nach einem weiteren Versuch gelegen haben. Die Wiederholungsrate darf deshalb nicht unmittelbar als Beleg für eine problematische Kameraführung interpretiert werden.

\begin{table}[htbp]
\centering
\small
\begin{tabular}{p{0.46\textwidth} c c c}
\hline
\textbf{Aussage} & \textbf{M} & \textbf{Md} & \textbf{Zust.} \\ \hline
Die Gesamtbewertung war für mich nachvollziehbar & 4,11 & 4 & 95\,\% \\
Ich habe verstanden, was die einzelnen Kriterien bedeuten & 4,37 & 4 & 89\,\% \\
Die Anzeige mit Prozentwerten und Sternen war verständlich & 4,58 & 5 & 89\,\% \\
Es war schwer zu erkennen, warum ich diese Bewertung bekommen habe ($-$) & 2,58 & 2 & 26\,\% \\ \hline
\end{tabular}
\caption{Aussagen zur Verständlichkeit der Bewertung ($n=\Nteiln$)}
\label{tab:eval_verstaendlich}
\end{table}

Die Ergebnisse zur Verständlichkeit fallen insgesamt positiv aus. 95 Prozent hielten die Gesamtbewertung für nachvollziehbar und 89 Prozent gaben an, die einzelnen Kriterien verstanden zu haben. H3 wird damit gestützt.

Gleichzeitig stimmten fünf Personen der Aussage zu, dass es schwer gewesen sei zu erkennen, warum sie genau diese Bewertung erhalten hatten. Dies deutet auf zwei unterschiedliche Ebenen von Nachvollziehbarkeit hin. Ein Ergebnis kann inhaltlich plausibel und die Bedeutung eines Kriteriums verständlich sein, ohne dass der konkrete Weg vom Messwert zur Punktzahl erkennbar ist. Die offenen Rückmeldungen in Abschnitt~\ref{sec:eval_freitext} bestätigen diese Unterscheidung. Die Nutzerstudie prüft damit die wahrgenommene Verständlichkeit der Darstellung, nicht die technische Rückverfolgbarkeit der Berechnung selbst.

\begin{table}[htbp]
\centering
\small
\begin{tabular}{p{0.46\textwidth} c c c}
\hline
\textbf{Aussage} & \textbf{M} & \textbf{Md} & \textbf{Zust.} \\ \hline
Ich vertraue der Bewertung der App & 4,16 & 4 & 89\,\% \\
Die Bewertung passte zu meiner eigenen Einschätzung & 4,00 & 4 & 79\,\% \\
Die App bewertet meinen Kick teilweise falsch ($-$) & 2,58 & 2 & 21\,\% \\ \hline
\end{tabular}
\caption{Aussagen zum Vertrauen in die Bewertung ($n=\Nteiln$)}
\label{tab:eval_vertrauen}
\end{table}

Auch das Vertrauen in die Bewertung fällt überwiegend positiv aus. 89 Prozent gaben an, der Bewertung zu vertrauen. 79 Prozent sahen eine Übereinstimmung mit ihrer eigenen Einschätzung. Gleichzeitig stimmten 21 Prozent der Aussage zu, dass die App ihren Kick teilweise falsch bewertet habe. H4 wird deskriptiv gestützt, die Antworten zeigen jedoch deutlicher als bei den übrigen Bereichen Vorbehalte gegenüber der Richtigkeit einzelner Bewertungen.

Die Übereinstimmung mit der Selbsteinschätzung ist nicht als Nachweis der Messgenauigkeit zu verstehen. Wie zuverlässig die Teilnehmenden die eigene technische Ausführung beurteilen, wurde nicht erhoben. Ohne unabhängige Referenz lässt sich daher weder aus einer Übereinstimmung noch aus einer Abweichung ableiten, ob die Systembewertung fachlich richtig war.

\begin{table}[htbp]
\centering
\small
\begin{tabular}{p{0.46\textwidth} c c c}
\hline
\textbf{Aussage} & \textbf{M} & \textbf{Md} & \textbf{Zust.} \\ \hline
Die Tipps haben mir geholfen zu verstehen, was ich verbessern kann & 4,21 & 4 & 89\,\% \\
Die App könnte mir helfen, meine Technik zu verbessern & 4,63 & 5 & 100\,\% \\
Die App könnte mir helfen, auch ohne Trainer zu arbeiten & 4,37 & 4 & 95\,\% \\
Für meine Technikverbesserung bringt mir die App wenig ($-$) & 1,68 & 2 & 0\,\% \\ \hline
\end{tabular}
\caption{Aussagen zum wahrgenommenen Nutzen ($n=\Nteiln$)}
\label{tab:eval_nutzen}
\end{table}

Der wahrgenommene Nutzen erhält die stärkste Zustimmung der vier untersuchten Bereiche. Alle Teilnehmenden stimmten der Aussage zu, dass die Anwendung ihnen helfen könnte, ihre Technik zu verbessern. Niemand stimmte der negativ formulierten Aussage zu, dass die Anwendung für die Technikverbesserung wenig bringe. H5 wird damit gestützt. Diese Ergebnisse beschreiben jedoch eine Erwartung beziehungsweise Wahrnehmung der Teilnehmenden. Ob die Anwendung tatsächlich zu einer messbaren Technikverbesserung führt, wurde nicht untersucht.

Die Nutzungsabsicht H6 wurde als Mehrfachauswahl erhoben. Alle \Nteiln{} Personen gaben an, die Anwendung zur Vorbereitung auf Gürtelprüfungen einsetzen zu wollen. 18 Personen würden sie weiterempfehlen. H6 wird damit deskriptiv gestützt. Aufgrund des binären Antwortformats lässt sich allerdings nicht erfassen, wie stark diese Absicht ausgeprägt ist. Für eine erneute Erhebung wäre eine abgestufte Skala aussagekräftiger.

\subsection{Ergebnisse zu Erkennung, Vokabelfunktion und Bedienbarkeit}
\label{sec:eval_weitere}
Der Block zur technischen Erkennung erhebt, ob während der Nutzung Probleme bei der Personenerkennung wahrgenommen wurden, bei welchen Modellvarianten sie auftraten und wie oft eine Aufnahme wiederholt wurde. Diese Angaben beschreiben subjektiv beobachtete technische Probleme, stellen jedoch keine kontrollierte Messung der Modellgenauigkeit dar.

\begin{table}[htbp]
\centering
\small
\begin{tabular}{p{0.46\textwidth} p{0.38\textwidth}}
\hline
\textbf{Frage} & \textbf{Ergebnis} \\ \hline
Probleme bei der Personenerkennung & 6 ja, 13 nein \\
Betroffene Modellvarianten (Mehrfachnennung) & Heavy 4, Lite 2, Full 2, alle Varianten 1 \\
Notwendige Wiederholungen der Aufnahme & 6 gar nicht, 9 einmal, 4 zweimal \\ \hline
\end{tabular}
\caption{Angaben zur technischen Erkennung ($n=\Nteiln$)}
\label{tab:eval_erkennung}
\end{table}

Sechs der \Nteiln{} Teilnehmenden berichteten von Problemen bei der Personenerkennung. Vier nannten dabei die Variante \textit{Heavy}, jeweils zwei \textit{Lite} und \textit{Full}, eine Person gab Probleme bei allen Varianten an. Daraus lässt sich kein Vergleich der Modellgüte ableiten, da nicht erhoben wurde, wie häufig jede Variante insgesamt genutzt wurde und die Videos nicht kontrolliert mit allen drei Varianten ausgewertet wurden. Ein systematischer Modellvergleich müsste dieselben Videos unter identischen Bedingungen mit jeder Variante analysieren.

Eine Antwort ist in sich widersprüchlich. Eine Person verneinte zunächst Erkennungsprobleme, wählte in der Folgefrage aber dennoch \textit{Heavy} aus und gab zwei Wiederholungen an. Möglicherweise wurde die Wiederholung nicht als Erkennungsproblem interpretiert. Die Antwort zeigt zugleich, dass die Fragen zur technischen Erkennung in einer Folgestudie präziser voneinander getrennt werden sollten.

Die Vokabelfunktion erreicht sehr hohe Zustimmungswerte. Die Verständlichkeit der Gurtauswahl und die Passung der angezeigten Vokabeln liegen jeweils bei einem Mittelwert von 4,89. Der wahrgenommene Nutzen für die Prüfungsvorbereitung erreicht 4,79. Alle drei Aussagen erhielten von allen Teilnehmenden Zustimmung.

Diese Ergebnisse sind nur ergänzend zu betrachten, da die Vokabelfunktion nicht Gegenstand der Forschungsfragen ist. Bei der bevorzugten Lernform nannten zwölf Personen Karteikarten, neun die Auswahl aus Vorgaben und sechs eine freie Texteingabe. Zwei Personen gaben keine Präferenz an. Die im Prototyp umgesetzte Multiple-Choice Abfrage entspricht damit nicht der am häufigsten genannten Präferenz.

Auch die Gestaltung wurde positiv bewertet. Navigation und Übersichtlichkeit erreichen jeweils 4,58, die Lesbarkeit von Schriftgrößen und Farben 4,74. Der allgemeine Designeindruck liegt bei 4,32. Die offenen Rückmeldungen zeigen dennoch konkrete Verbesserungsmöglichkeiten, insbesondere bei der Auffindbarkeit des Uploads und beim Debug-Overlay.

Die Gebrauchstauglichkeit wurde zusätzlich mit der \textit{System Usability Scale} erhoben \cite{Brooke1996}. Bei den positiv formulierten Items wird jeweils 1 vom Skalenwert abgezogen. Bei den negativ formulierten Items wird der Skalenwert von 5 subtrahiert. Die Summe der zehn Beiträge wird mit 2,5 multipliziert und ergibt einen SUS-Wert zwischen 0 und 100.

Die Anwendung erreicht einen Mittelwert von \SUSm{} bei einer Standardabweichung von \SUSsd{}. Der Median liegt bei \SUSmd{}, die Einzelwerte reichen von \SUSmin{} bis \SUSmax{}. Als Vergleichsmaßstab wird der häufig verwendete durchschnittliche SUS-Wert von 68 herangezogen \cite{SauroLewis2016}. Siebzehn der \Nteiln{} Teilnehmenden liegen über diesem Wert, sechzehn erreichen mehr als 80 Punkte. Der niedrigste Einzelwert beträgt \SUSmin{} und liegt damit knapp unter dem Vergleichswert. H1 wird damit gestützt.

\subsection{Offene Rückmeldungen}
\label{sec:eval_freitext}
Von den \Nteiln{} Teilnehmenden beantworteten \Nfreitext{} mindestens eine der offenen Fragen. Die Antworten wurden nach wiederkehrenden Themen gruppiert und zusammenfassend wiedergegeben.

\subsubsection{Positiv hervorgehobene Eigenschaften}
Am häufigsten wurde die Videoanalyse selbst positiv hervorgehoben. Eine Person betonte, dass die Anwendung Verbesserungsvorschläge in natürlicher Sprache liefere und nicht nur einen Punktwert anzeige. Ebenfalls positiv wurde das annotierte Video mit dem erkannten Skelett bewertet. Weitere Antworten hoben hervor, dass eigene Fehler verständlich benannt würden. Diese Rückmeldungen stützen die Entscheidung, bei schwach bewerteten Kriterien neben dem Prozentwert auch einen sprachlichen Verbesserungshinweis anzuzeigen.

Mehrfach positiv genannt wurden außerdem die Bedienbarkeit, der schnelle Videoupload, die Speicherung früherer Analysen im Verlauf und die Vokabelfunktion. Eine Person beschrieb zudem einen Motivationseffekt der Punktebewertung, da der Versuch, einen höheren Wert zu erreichen, einen spielerischen Anreiz erzeuge.

\subsubsection{Als unklar oder störend genannte Punkte}
Mehrere Rückmeldungen beziehen sich unmittelbar auf Entscheidungen des Prototyps. Eine Person gab an, den Begriff \textit{Chamber} nicht verstanden zu haben. Der Begriff wird in der Benutzeroberfläche verwendet, dort jedoch nicht erklärt.

Drei Personen sprachen die Auswahl zwischen \textit{Lite}, \textit{Full} und \textit{Heavy} an. Unter anderem wurde gefragt, weshalb eine kleinere Variante genutzt werden sollte, wenn \textit{Heavy} zur Verfügung steht, und es wurde eine Erklärung der Unterschiede gewünscht. Dies bestätigt, dass die Modellauswahl eher eine Entwicklungsfunktion als eine für Endnutzende relevante Entscheidung darstellt.

Das Debug-Overlay im Ausgabevideo wurde ebenfalls als störend beschrieben. Eine Person empfand die schnell wechselnden Werte als verwirrend. Da das Overlay in erster Linie der Entwicklung und Fehlersuche dient, sollte es in einer Endanwendung standardmäßig ausgeblendet werden.

Zwei Personen fanden die Upload-Schaltfläche erst nach dem Scrollen. Da der Upload den Einstieg in die Analyse bildet, ist dies trotz der einfachen gestalterischen Ursache ein relevanter Usability-Punkt.

Die inhaltlich wichtigste konkrete Freitext-Rückmeldung betrifft die wahrgenommene Richtigkeit einzelner Kriterienwerte. Eine Person berichtete, dass die Hüftausrichtung beziehungsweise Rotation teilweise als zu stark bewertet worden sei und das Kriterium 0 Prozent erhalten habe. Außerdem sei die Beinstreckung teilweise niedrig bewertet worden. Beim \textit{Roundhouse Kick} habe die Anwendung mitunter ein Absinken des Knies um 30 bis 50 Zentimeter angezeigt, obwohl dies aus Sicht der Person nicht zutraf. Diese Aussage ist die einzige konkrete Freitext-Rückmeldung zu möglichen Mess- beziehungsweise Bewertungsfehlern. Zusammen mit der quantitativen Aussage, dass 21 Prozent zumindest teilweise eine falsche Bewertung wahrnahmen, zeigt sie, welche Größen in einer späteren Validierung besonders geprüft werden sollten. Ohne unabhängige Referenz bleibt offen, ob die Abweichung durch das System, die Aufnahmebedingungen oder die Selbsteinschätzung entstand.

\subsubsection{Gewünschte Funktionen}
Am deutlichsten formuliert wurde der Wunsch nach einer Erklärung der Punktberechnung. Als Beispiel wurde genannt, dass bei einem Standbeinwert von 60 Prozent sichtbar sein sollte, weshalb genau dieser Wert zustande kommt und was für einen höheren Wert verändert werden müsste.

Die dafür erforderlichen Informationen liegen grundsätzlich im System vor. Für jedes Kriterium existieren eine Messgröße sowie die beiden Grenzwerte der Bewertungsfunktion. Diese Informationen werden in der Benutzeroberfläche bislang jedoch nicht vollständig dargestellt. Die technische Rückverfolgbarkeit der Bewertung ist damit vorhanden, wird aber nicht konsequent in eine für Nutzende verständliche Erklärung übersetzt.

Weitere Wünsche betreffen die direkte Videoaufnahme in der Anwendung, Beispielvideos zur korrekten Ausführung, zusätzliche Kick- und Schlagtechniken, Poomsae und Übungen. Für die Vokabelfunktion wurden unter anderem koreanische Schreibweisen, Aussprache, Karteikarten, freie Texteingabe, eine stärkere Auswahl nach Gurtgrad und eine Auflösung am Ende der Abfrage gewünscht.

\subsubsection{Anmerkungen zum Untersuchungsaufbau}
Eine Rückmeldung bezog sich auf die Notwendigkeit einer zweiten Person zum Filmen. Dies entspricht der Begründung für die Zweiergruppen, zeigt aber zugleich eine Einschränkung des Einsatzszenarios. Wer allein trainiert, benötigt für die aktuelle Form der Anwendung eine feste Kameraposition, beispielsweise über ein Stativ, oder eine weitere Person zum Filmen.

\subsection{Zusammenfassende Einordnung}
\label{sec:eval_einordnung}
Die sechs Hypothesen werden auf Grundlage der deskriptiven Ergebnisse gestützt. Tabelle~\ref{tab:eval_hypergebnis} fasst die Befunde zusammen.

\begin{table}[htbp]
\centering
\small
\begin{tabular}{p{0.05\textwidth} p{0.39\textwidth} p{0.24\textwidth} p{0.14\textwidth}}
\hline
& \textbf{Aussage} & \textbf{Ergebnis} & \textbf{Befund} \\ \hline
H1 & SUS über Vergleichswert 68 & \SUSm{}; 17 von \Nteiln{} darüber & gestützt \\
H2 & Aufnahme einfach und verständlich & Zustimmung 79 bis 100\,\% in positiver Richtung & gestützt \\
H3 & Bewertung verständlich und nachvollziehbar & Zustimmung 74 bis 95\,\% in positiver Richtung & gestützt \\
H4 & Vertrauen in die Bewertung & Zustimmung 79 bis 89\,\% bei positiven Aussagen & gestützt \\
H5 & Wahrgenommener Nutzen & Zustimmung 89 bis 100\,\% & gestützt \\
H6 & Nutzungsabsicht & 19 Prüfungsvorbereitung, 18 Weiterempfehlung & gestützt \\ \hline
\end{tabular}
\caption{Deskriptive Einordnung der Hypothesen ($n=\Nteiln$)}
\label{tab:eval_hypergebnis}
\end{table}

Die überwiegend positiven Werte müssen im Kontext des Untersuchungsaufbaus interpretiert werden. Die Stichprobe ist klein, nicht zufällig ausgewählt und überwiegend aus unteren Gurtgraden zusammengesetzt. Die Anwendung war für die Teilnehmenden neu und es gab keine Vergleichsanwendung. Zudem probierten die Personen den Prototyp in Zweiergruppen ohne standardisierte Aufgabenfolge oder einheitliche Nutzungsdauer aus. Die Ergebnisse lassen sich daher nicht ohne Weiteres auf eine größere Grundgesamtheit übertragen.

Inhaltlich ergeben sich drei zentrale Befunde. Erstens wird die regelbasierte Rückmeldung von der Zielgruppe überwiegend verstanden und als hilfreich wahrgenommen. Zweitens reicht diese Verständlichkeit nicht bei allen Personen bis zur Erklärung des konkreten Punktwerts. Die Anwendung zeigt Kriterien und Prozentwerte, macht den Weg vom Messwert über die Grenzwerte zur Punktzahl aber nicht sichtbar. Drittens wurden trotz positiv bewerteter Aufnahmesituation häufig Videos wiederholt. Die vorhandenen Daten erlauben jedoch keine eindeutige Zuordnung dieser Wiederholungen zu Aufnahme, Erkennung oder Ausführung.

Die Nutzerstudie beantwortet damit vor allem die zweite Forschungsfrage. Sie zeigt eine hohe wahrgenommene Gebrauchstauglichkeit und überwiegend positive Einschätzungen von Verständlichkeit, Vertrauen und Nutzen. Sie beantwortet nicht, wie genau die zugrunde liegenden Messgrößen sind oder ob die Punktwerte mit fachlichen Trainerurteilen übereinstimmen. Diese offene Validierungsfrage ist insbesondere für den \textit{Spinning Back Kick} relevant, der nur von 14 Personen getestet wurde und zugleich die größten konzeptionellen Herausforderungen aufweist.

\section{Fazit \& Ausblick}

\subsection{Zusammenfassung}
Die vorliegende Arbeit hat einen webbasierten Prototyp entwickelt, der einzelne, mit einem Smartphone aufgenommene Taekwondo-Kicks mithilfe markerloser \textit{Pose Estimation} analysiert. Aus den von MediaPipe gelieferten Landmarken werden ausgewählte Winkel, Positionen und Richtungsgrößen berechnet. Eine technikspezifische Zustandsmaschine ordnet die Bewegung verschiedenen Phasen zu. Anschließend werden sieben Kriterien für den \textit{Front Kick} und den \textit{Roundhouse Kick} sowie fünf Kriterien für den \textit{Spinning Back Kick} regelbasiert bewertet.

Die erste Forschungsfrage betrifft die prototypische technische Umsetzung und deren Grenzen. Die Arbeit zeigt, dass sich die ausgewählten Merkmale aus acht Körperlandmarken ableiten und in einem einheitlichen Bewertungsschema verarbeiten lassen. Anders als bei Wang et al. \cite{Wang2026} wird die Technik vor der Analyse vom Nutzer ausgewählt. Dadurch ist kein zusätzlich trainiertes Modell zur Technikerkennung erforderlich. Die eigentliche Pose Estimation basiert weiterhin auf den trainierten Modellen von MediaPipe. Die Bewertung selbst erfolgt jedoch anhand explizit definierter Messgrößen und Grenzwerte.

Die Umsetzung zeigt zugleich, dass die Nachvollziehbarkeit eines regelbasierten Verfahrens nicht automatisch fachliche Validität bedeutet. Die Grenzwerte wurden nur teilweise aus publizierter Kinematik abgeleitet und im Übrigen aus Trainingspraxis und prototypischer Entwicklung festgelegt. Eine unabhängige messtechnische Validierung und eine Kalibrierung gegenüber Trainerurteilen stehen aus. Die Arbeit kann daher zeigen, wie die Kriterien berechnet werden und an welchen Stellen technische Grenzen auftreten, nicht jedoch, wie genau die erzeugten Punktwerte die tatsächliche Qualität eines Kicks abbilden.

Die zweite Forschungsfrage wurde mit einer Nutzerstudie mit \Nteiln{} Taekwondo-Trainierenden untersucht. Die Anwendung erreicht auf der \textit{System Usability Scale} einen Mittelwert von \SUSm{} Punkten. Die Teilnehmenden bewerteten die Verständlichkeit und den wahrgenommenen Nutzen überwiegend positiv. Alle 19 Personen gaben an, die Anwendung zur Prüfungsvorbereitung einsetzen zu wollen, 18 würden sie weiterempfehlen. Besonders positiv wurde hervorgehoben, dass schwache Kriterien nicht nur als Prozentwert angezeigt, sondern mit einem sprachlichen Verbesserungshinweis verbunden werden.

Gleichzeitig zeigte die Studie eine Lücke zwischen verständlicher Rückmeldung und erklärter Bewertung. Die Anwendung nennt das jeweilige Kriterium und den erreichten Prozentwert, zeigt aber nicht, wie der konkrete Messwert im Verhältnis zu den Grenzwerten zur Punktzahl geführt hat. Diese Information ist im regelbasierten System grundsätzlich vorhanden, wird in der Benutzeroberfläche jedoch noch nicht vollständig vermittelt.

\subsection{Diskussion der Grenzen}
Die Grenzen der Arbeit lassen sich in messtechnische, konzeptionelle und methodische Aspekte unterteilen.

\subsubsection{Messtechnische Grenzen}
Eine zentrale Grenze betrifft die körperrelative Hüftausrichtung. Sie wird aus den geschätzten 3D Positionen der beiden Hüften und der Richtung des Kickbeins abgeleitet (Abschnitt~4.7). Bei seitlicher Körperausrichtung liegen die projizierten Hüftlandmarken im Kamerabild näher beieinander und können sich teilweise verdecken. MediaPipe muss die Position eines verdeckten Punkts dann stärker schätzen. Dadurch kann die abgeleitete Beckenrichtung gerade bei rotationsstarken Bewegungen unsicher werden.

Die aktuelle Implementierung nutzt die von MediaPipe gelieferten Sichtbarkeitswerte nicht, um unsichere Landmarken zu erkennen oder Messwerte zu verwerfen. Der vorhandene Sprungfilter reduziert zwar abrupte Änderungen der Hüftausrichtung, kann aber keine kleineren, plausibel wirkenden Schätzfehler erkennen.

Mehrere Kriterien beruhen außerdem auf einzelnen Extremwerten innerhalb einer Phase. Beinstreckung und Chamber-Winkel verwenden jeweils einen maximalen beziehungsweise minimalen Wert. Die Ausrichtung im Treffmoment wird aus einem einzelnen Referenzframe übernommen. Beim Knieabsinken wird die Differenz zwischen zwei Extremwerten verwendet. Ein einzelner ungenauer Landmark kann solche Werte deshalb stärker beeinflussen als ein über mehrere Bilder aggregiertes Maß. Beim Standbeinwinkel wird bereits ein Median verwendet. Ein ähnlicher Ansatz über Median, Perzentile oder kurze Zeitfenster um den Treffmoment könnte auch andere Kriterien robuster machen.

Eine weitere Grenze betrifft die Geschwindigkeitsanzeige. Die Weltkoordinaten von MediaPipe sind auf die Hüftmitte bezogen. Die berechnete Geschwindigkeit beschreibt deshalb die Bewegung des Fußes relativ zur Hüfte und nicht die vollständige Bewegung des Fußes im Raum. Eine globale Translation des Körpers wird nicht berücksichtigt. Da die Geschwindigkeit nicht in die Bewertung eingeht, betrifft diese Einschränkung nur die Zusatzanzeige.

Schließlich wird das Knieabsinken in absoluten Metern gemessen, während längenbasierte Größen beim \textit{Spinning Back Kick} auf die Beinlänge normiert werden. Dadurch ist die Bewertung dieses Kriteriums nicht vollständig körpergrößenunabhängig.

\subsubsection{Konzeptionelle Grenzen}
Die Zustandsmaschinen hängen von festen Übergangsschwellen ab. Ein Teil dieser Werte orientiert sich an den Bewertungskriterien, andere wurden heuristisch anhand der Bewegungsstruktur und von Testaufnahmen festgelegt. Diese Schwellen wurden nicht anhand eines annotierten Datensatzes optimiert. Fehler in der Phasenerkennung können sich deshalb auf alle anschließend aus der Phase berechneten Kriterien auswirken.

Wird die Kick-Phase überhaupt nicht erkannt, bricht die Auswertung früh ab und gibt nur die Beinstreckung mit 0 Prozent zurück (Abschnitt~4.9). Eine Ausführung knapp unterhalb der Erkennungsschwelle erhält damit dasselbe Gesamtergebnis wie ein Video ohne verwertbaren Kick. In einer Endanwendung sollte dieser Zustand deshalb nicht als schlechte technische Bewertung, sondern als fehlgeschlagene Analyse kommuniziert werden.

Die Gesamtbewertung ist ein ungewichteter Mittelwert. Dadurch zählt jedes Kriterium gleich stark, obwohl einzelne Merkmale für die Identität und Qualität einer Technik unterschiedlich bedeutsam sein können. Eine belastbare Gewichtung würde jedoch eine fachliche Referenz voraussetzen und wurde deshalb nicht nachträglich ohne empirische Grundlage eingeführt.

Auch die Farbgebung der Kriterien ist nicht vollständig mit der Bestehenslogik abgestimmt. Ein Wert zwischen 50 und 74 Prozent gilt als bestanden und erzeugt keinen Verbesserungstipp, wird in der aktuellen Benutzeroberfläche aber rot dargestellt. Diese Inkonsistenz sollte in einer Weiterentwicklung behoben werden.

Der Prototyp verarbeitet nur einen Kick pro Video. Die Kickseite wird einmal festgelegt und anschließend nicht zurückgesetzt. Serien mehrerer Kicks oder ein Wechsel der Kickseite innerhalb einer Aufnahme können daher nicht ausgewertet werden.

\subsubsection{Besonderheiten beim Spinning Back Kick}
\label{sec:grenzen_dwit}
Der \textit{Spinning Back Kick} bündelt mehrere der genannten Grenzen. Wang et al. berichten bereits beim \textit{Roundhouse Kick} von Problemen durch Rotation und Selbstverdeckung \cite{Wang2026}. Bei einer etwa 180 Grad umfassenden Körperdrehung treten diese Herausforderungen noch ausgeprägter auf.

Erstens bildet die aufsummierte Beckendrehung nicht den vollständigen Rotationsweg ab. Die Kickseite wird erst festgelegt, wenn eine ausreichende Höhendifferenz zwischen den Knien erkannt wurde. Frühe Anteile der Eindrehung können dadurch fehlen. Zusätzlich werden Winkeländerungen über 30 Grad pro Bild als unplausibel verworfen. Bei schnellen Bewegungen kann dies reale Rotationsanteile entfernen.

Zweitens misst die im Treffmoment verwendete körperrelative Hüftausrichtung keine globale Körperdrehung. Sie beschreibt die Lage des Kickbeins relativ zum Becken. Eine Person kann sich weit im Raum eingedreht haben und dennoch einen ungünstigen körperrelativen Winkel erreichen. Die Rückmeldung muss deshalb klar zwischen globaler Eindrehung und relativer Ausrichtung von Becken und Kickbein unterscheiden.

Drittens ist das Winkelmaß auf 180 Grad begrenzt. Es kann nicht unterscheiden, ob eine Ausrichtung den Grenzbereich gerade erreicht oder bereits darüber hinaus weitergedreht wurde. Die zusätzlich geführte Nachdrehung dient deshalb lediglich zur Auswahl eines Rückmeldungstexts und geht nicht in die Punktzahl ein.

Viertens erhöht die Selbstverdeckung während der Drehung die Unsicherheit der Hüft, Knie und Sprunggelenklandmarken. Gerade die für das Kriterium zentrale Beckenrichtung kann dadurch in denjenigen Frames unsicher werden, in denen die Rotation besonders stark ausgeprägt ist.

Die erste Forschungsfrage lässt sich damit vorsichtig beantworten. Für alle drei Techniken konnten die ausgewählten Messgrößen in einem regelbasierten Prototyp umgesetzt werden. Beim \textit{Spinning Back Kick} zeigen sich jedoch deutlich stärkere Grenzen durch Selbstverdeckung, heuristische Rotationszählung und die körperrelative Definition des verwendeten Winkelmaßes. Ob diese Effekte lediglich die Robustheit reduzieren oder zu systematischen fachlichen Fehlbewertungen führen, kann ohne unabhängige Referenz nicht entschieden werden.

\subsubsection{Grenzen der Nutzerstudie}
Auch die Nutzerstudie besitzt Einschränkungen. Die Stichprobe umfasst nur 19 Personen und ist eine Gelegenheitsstichprobe. Die unteren Gurtgrade sind deutlich stärker vertreten als fortgeschrittene Trainierende. Die Ergebnisse beschreiben daher vor allem die Wahrnehmung dieser konkreten Gruppe.

Die Anwendung wurde in Zweiergruppen frei ausprobiert. Es gab keine standardisierte Aufgabenfolge, keine einheitliche Nutzungsdauer und keine Vergleichsanwendung. Die Personen konnten sich während der Erprobung gegenseitig beeinflussen. Zudem beruhen Aussagen zu Vertrauen und wahrgenommener Richtigkeit auf subjektiven Einschätzungen. Sie ersetzen weder eine technische Referenzmessung noch ein unabhängiges fachliches Urteil.

Die Nutzungsabsicht wurde lediglich über zwei Auswahloptionen erfasst. Deshalb lässt sich nicht bestimmen, wie stark die Bereitschaft zur späteren Nutzung tatsächlich ausgeprägt ist. Auch die häufigen Wiederholungen von Videoaufnahmen können nicht eindeutig einer Ursache zugeordnet werden, da der Fragebogen nicht zwischen Aufnahmefehler, Erkennungsproblem, misslungener Ausführung und freiwilligem weiteren Versuch unterscheidet.

\subsection{Ausblick}
\label{sec:ausblick}
Aus den technischen und nutzerbezogenen Ergebnissen ergeben sich mehrere Weiterentwicklungen.

Kurzfristig sollte zunächst die Rückmeldung bei nicht erkannter Kick-Phase geändert werden. Statt eines Gesamtwerts von 0 Prozent sollte die Anwendung deutlich anzeigen, dass keine verwertbare Kick-Phase erkannt wurde. Ebenso sollten die Farbgrenzen der Kriterien an die Bestehensgrenze angepasst werden. Die Modellauswahl und das Debug-Overlay sollten für reguläre Nutzende ausgeblendet werden. Zusätzlich sollte die Benutzeroberfläche erklären, wie ein Prozentwert aus Messwert und Grenzwerten entsteht.

Für die Robustheit der Messgrößen bieten sich mehrere technische Änderungen an. Sichtbarkeitswerte von MediaPipe könnten verwendet werden, um unsichere Landmarken zu erkennen. Einzelne Extremwerte könnten durch Perzentile, Mediane oder kurze Zeitfenster ersetzt werden. Das Knieabsinken sollte wie die längenbasierten Größen des \textit{Spinning Back Kick} auf die Beinlänge normiert werden. Für die Beckenrichtung könnte zusätzlich geprüft werden, ob die zugrunde liegenden Hüftlandmarken im jeweiligen Frame ausreichend zuverlässig sind.

Mittelfristig sollten weitere Landmarken einbezogen werden. Ferse und Zehen würden es ermöglichen, die Ausrichtung des Standfußes und die Trefferfläche beziehungsweise Fußstellung zu berücksichtigen. Damit ließen sich beispielsweise der Pivot des Standfußes beim \textit{Roundhouse Kick} und zusätzliche Ausrichtungsmerkmale beim \textit{Spinning Back Kick} untersuchen.

Der wichtigste nächste Schritt ist jedoch eine zweistufige Validierung. Zunächst sollte die technische Messgüte geprüft werden. Dazu können ausgewählte Frames manuell annotiert oder mit einem geeigneten Referenzsystem verglichen werden. So ließe sich untersuchen, wie stark Kniewinkel, Hüftwinkel und körperrelative Hüftausrichtung von einer Referenz abweichen und wie sich Kameraansicht, Sichtbarkeit und Modellvariante auf diese Fehler auswirken.

Davon getrennt sollte die fachliche Bewertungsvalidität untersucht werden. Erfahrene Trainerinnen und Trainer könnten dieselben Videos unabhängig beurteilen. Anschließend ließe sich prüfen, ob die Reihenfolge oder kategoriale Einordnung der Systembewertungen mit den Trainerurteilen übereinstimmt. Erst dieser Vergleich kann zeigen, ob die gewählten Grenzwerte und der ungewichtete Gesamtwert technisch plausible Messgrößen auch in eine fachlich sinnvolle Bewertung überführen.

Für den \textit{Spinning Back Kick} sollte zusätzlich eine robustere Repräsentation der globalen Körperrotation untersucht werden. Die aktuelle körperrelative Hüftausrichtung und die aufsummierte Änderung der Beckenrichtung erfassen unterschiedliche Aspekte der Bewegung und besitzen jeweils eigene Grenzen. Eine Kombination zusätzlicher Landmarken, zeitlicher Rotationsmerkmale und einer expliziten Behandlung verdeckter Körperpunkte könnte hier zuverlässiger sein.

Erst mit einer solchen technischen und fachlichen Validierung lässt sich abschließend beurteilen, wie genau die berechneten Bewertungen die tatsächliche technische Qualität der drei Taekwondo-Kicks abbilden und ob rotationsstarke Techniken mit demselben regelbasierten Grundansatz zuverlässig bewertet werden können.

\appendix
\includepdf[
    pages=-,
    fitpaper=true
]{fragebogen.pdf}
\begin{thebibliography}{00}
% [1]
\bibitem{Colyer2018}
Colyer, S.~L., Evans, M., Cosker, D.~P. \& Salo, A.~I.~T.
\textit{A Review of the Evolution of Vision-Based Motion Analysis
and the Integration of Advanced Computer Vision Methods
Towards Developing a Markerless System}.
Sports Medicine -- Open, 4(1): 24, 2018.
\href{https://doi.org/10.1186/s40798-018-0139-y}{https://doi.org/10.1186/s40798-018-0139-y}

% [2]
\bibitem{Edriss2025}
Edriss, S., Romagnoli, C., Caprioli, L., Bonaiuto, V.,
Padua, E. \& Annino, G.
\textit{Commercial vision sensors and AI-based pose
estimation frameworks for markerless motion analysis
in sports and exercises: a mini review}.
Frontiers in Physiology, 16: 1649330, 2025.
\href{https://doi.org/10.3389/fphys.2025.1649330}{https://doi.org/10.3389/fphys.2025.1649330}

% [3]
\bibitem{Babouras2024}
Babouras, A., Abdelnour, P., Fevens, T. \& Martineau, P.~A.
\textit{Comparing novel smartphone pose estimation frameworks
with the Kinect V2 for knee tracking during athletic stress tests}.
International Journal of Computer Assisted Radiology and Surgery,
19(7): 1321--1328, 2024.
\href{https://doi.org/10.1007/s11548-024-03156-5}{https://doi.org/10.1007/s11548-024-03156-5}

% [4]
\bibitem{YilmazAtes2025}
Aybüke Yılmaz, E. \& Ateş, O.
\textit{Taekwondo kicks at the center of biomechanical
research: A systematic review}.
Kinesiology, 57(2): 289--307, 2025.
\href{https://doi.org/10.26582/k.57.2.13}{https://doi.org/10.26582/k.57.2.13}

% [5]
\bibitem{Falco2009}
Falco, C., Alvarez, O., Castillo, I., Estevan, I., Martos, J.,
Mugarra, F. \& Iradi, A.
\textit{Influence of the distance in a roundhouse kick's
execution time and impact force in Taekwondo}.
Journal of Biomechanics, 42(3): 242--248, 2009.
\href{https://doi.org/10.1016/j.jbiomech.2008.10.041}{https://doi.org/10.1016/j.jbiomech.2008.10.041}

% [6]
\bibitem{Wang2026}
Wang, Z., Sun, S., Liang, X. \& Chen, L.
\textit{Recognition and explainable quantitative evaluation of
fundamental taekwondo kicks from smartphone videos}.
Sensors, 26(14): 4499, 2026.
\href{https://doi.org/10.3390/s26144499}{https://doi.org/10.3390/s26144499}

% [7]
\bibitem{Moedinger2022}
M.~Mödinger, A.~Woll und I.~Wagner,
\textit{Video-based visual feedback to enhance motor learning in physical education -- a systematic review},
German Journal of Exercise and Sport Research, Bd. 52, S. 447--460, 2022.
doi:10.1007/s12662-021-00782-y.

% [8]
\bibitem{Pham2022}
Pham, Q.-T., Nguyen, V.-A., Nguyen, T.-T., Nguyen, D.-A.,
Nguyen, D.-G. \& Pham, D.-T.
\textit{Automatic recognition and assessment of physical
exercises from RGB images}.
In: 2022 IEEE Ninth International Conference on Communications
and Electronics (ICCE), Nha Trang, Vietnam, 349--354, 2022.
\href{https://doi.org/10.1109/ICCE55644.2022.9852094}{https://doi.org/10.1109/ICCE55644.2022.9852094}

% [9]
\bibitem{ElRajab2025}
El-Rajab, I., Klotzbier, T.~J., Korbus, H. \& Schott, N.
\textit{Camera-based mobile applications for movement
screening in healthy adults: a systematic review}.
Frontiers in Sports and Active Living, 7: 1531050, 2025.
\href{https://doi.org/10.3389/fspor.2025.1531050}{https://doi.org/10.3389/fspor.2025.1531050}

% [10]
\bibitem{Fong2011}
Fong, S.~S.~M. \& Ng, G.~Y.~F.
\textit{Does Taekwondo training improve physical fitness?}
Physical Therapy in Sport, 12(2): 100--106, 2011.
\href{https://doi.org/10.1016/j.ptsp.2010.07.001}{https://doi.org/10.1016/j.ptsp.2010.07.001}

% [11]
\bibitem{WTKyorugi2026}
World Taekwondo,
\textit{Competition Rules \& Interpretation, in force as of June 1, 2026},
Seoul, 2026.
\url{https://www.worldtaekwondo.org/}

% [12]
\bibitem{WTPoomsae2024}
World Taekwondo,
\textit{Poomsae Competition Rules \& Interpretation,
in force as of September 30, 2024},
Seoul, 2024.
\url{https://www.worldtaekwondo.org/att\_file/documents/Poomsae\_Competition\_Rules\_and\_Interpretation\_(In\_force\_as\_of\_September\_30\_2024).pdf}

% [13]
\bibitem{WikiBild2026}
Wikipedia,
\textit{Taekwondo-Wettkampf Olympische Sommerspiele 2016},
Online verfügbar unter:
\url{https://de.wikipedia.org/wiki/Taekwondo#/media/Datei\:Milad\_Kharchegani\_at\_the\_2016\_Summer\_Olympics.jpg}
Lizenz: CC BY 4.0,
abgerufen am 19. März 2026.

% [14]
\bibitem{KukkiwonRules2018}
Kukkiwon,
\textit{The Rules of Taekwondo Promotion Test},
Seoul, 2018.
\url{https://www.kukkiwon.or.kr/eng/board/read?boardManagementNo=51&boardNo=1398&menuLevel=2&menuNo=97}

% [15]
\bibitem{wasik2016}
Wąsik, J.; Góra, T.
\textit{The Kinematics of Taekwon-Do Back Kick}.
\textit{Baltic Journal of Health and Physical Activity} \textbf{2016}, 8(4), 49--55.
\href{https://doi.org/10.29359/BJHPA.08.4.06}{https://doi.org/10.29359/BJHPA.08.4.06}

% [16]
\bibitem{Gavagan2017}
Gavagan, C.~J. \& Sayers, M.~G.~L.
\textit{A biomechanical analysis of the roundhouse kicking technique
of expert practitioners: A comparison between the martial arts
disciplines of Muay Thai, Karate, and Taekwondo}.
PLoS One, 12(8): e0182645, 2017.
\href{https://doi.org/10.1371/journal.pone.0182645}{https://doi.org/10.1371/journal.pone.0182645}

% [17]
\bibitem{Gulia2022}
S.~Gulia und D.~Shaw,
\textit{Two-dimensional biomechanical analysis of front kick and front toe kick of female national level taekwondo players},
International Journal of Physical Education, Sports and Health, Bd. 9, Nr. 5, S. 130--134, 2022.
\href{https://www.semanticscholar.org/paper/Two-dimensional-biomechanical-analysis-of-front-and-Gulia-Scholar/5c5a35662600f7d4c74289a1c994e5161f38b25e\#extracted}{https://www.semanticscholar.org/paper/Two-dimensional-biomechanical-analysis-of-front-and-Gulia-Scholar/5c5a35662600f7d4c74289a1c994e5161f38b25e\#extracted}

% [18]
\bibitem{TaekwonWikiaDol}
TaeKwonWikia,
\textit{Dollyeo-chagi},
Online verfügbar unter:
\url{https://taekwonwikia.com.br/dollyeo-chagi/}
abgerufen am 16.08.2026.

% [19]
\bibitem{pedzich2006}
Pędzich, W.; Mastalerz, A.; Urbanik, C.
\textit{The Comparison of the Dynamics of Selected Leg Strokes in Taekwondo WTF}.
\textit{Acta of Bioengineering and Biomechanics} \textbf{2006}, 8(1), 83--90.
\href{https://actabio.pwr.edu.pl/fcp/0GBUKOQtTKlQhbx08SlkTUARAUWRuHQwFDBoIVURNWHVaFVZpCFghUHcKVigEQUw/302/public/publikacje/v8-1-2006/8.pdf}{https://actabio.pwr.edu.pl/fcp/0GBUKOQtTKlQhbx08SlkTUARAUWRuHQwFDBoIVURNWHVaFVZpCFghUHcKVigEQUw/302/public/publikacje/v8-1-2006/8.pdf}

% [20]
\bibitem{TaekwonWikiaDwi}
TaeKwonWikia,
\textit{Dwi-chagi},
Online verfügbar unter:
\url{https://taekwonwikia.com.br/dwi-chagi/}
abgerufen am 16.08.2026.

% [21]
\bibitem{Thomas2026}
Thomas, A., Tucker, C.~B., Lunn, D.~E., Cooke, M.~J.,
Nicholson, G. \& Walker, J.
\textit{Lights, Cameras, Action -- The influence of lighting
and camera position on walking and running kinematic measurements
using Theia3D markerless motion capture}.
Journal of Biomechanics, 203: 113343, 2026.
\href{https://doi.org/10.1016/j.jbiomech.2026.113343}{https://doi.org/10.1016/j.jbiomech.2026.113343}

% [22]
\bibitem{MediaPipeTasks}
Google LLC,
\textit{MediaPipe Tasks Overview},
Online verfügbar unter:
\url{https://developers.google.com/mediapipe/solutions}
abgerufen am 16. März 2026.

% [23]
\bibitem{Lugaresi2019}
C.~Lugaresi, J.~Tang, H.~Nash, C.~McClanahan, E.~Uboweja, M.~Hays, F.~Zhang,
C.-L.~Chang, M.~Yong, J.~Lee, W.-T.~Chang, W.~Hua, M.~Georg und M.~Grundmann,
\textit{MediaPipe: A Framework for Perceiving and Processing Reality},
Third Workshop on Computer Vision for AR/VR at IEEE CVPR, 2019.
Online verfügbar unter:
\url{https://mixedreality.cs.cornell.edu/s/NewTitle\_May1\_MediaPipe\_CVPR\_CV4ARVR\_Workshop\_2019.pdf}

% [24]
\bibitem{Bazarevsky2020}
Bazarevsky, V., Grishchenko, I., Raveendran, K., Zhu, T., Zhang, F.
\& Grundmann, M.
\textit{BlazePose: On-device Real-time Body Pose Tracking}.
arXiv preprint arXiv:2006.10204, 2020.
\href{https://doi.org/10.48550/arXiv.2006.10204}{https://doi.org/10.48550/arXiv.2006.10204}

% [25]
\bibitem{mediapipepose}
Google LLC,
\textit{MediaPipe Pose Landmarks},
Online verfügbar unter:
\url{https://ai.google.dev/edge/mediapipe/solutions/vision/pose\_landmarker?hl=de}
abgerufen am 14. Mai 2026.

% [26]
\bibitem{mobilenetv2}
Sandler, M., Howard, A., Zhu, M., Zhmoginov, A., Chen, L.-C.
\textit{MobileNetV2: Inverted Residuals and Linear Bottlenecks}.
arXiv preprint arXiv:1801.04381,
2018.
\url{https://arxiv.org/abs/1801.04381}

% [27]
\bibitem{Tharatipyakul2024}
Tharatipyakul, A., Srikaewsiew, T. \& Pongnumkul, S.
\textit{Deep learning-based human body pose estimation in providing
feedback for physical movement: A review}.
Heliyon, 10(17): e36589, 2024.
\href{https://doi.org/10.1016/j.heliyon.2024.e36589}{https://doi.org/10.1016/j.heliyon.2024.e36589}

% [28]
\bibitem{Kim2011}
Kim, Y. K., Kim, Y. H., \& Im, S. J.
\textit{Inter-joint coordination in producing kicking velocity of taekwondo kicks}.
Journal of Sports Science and Medicine,
2011, 10(1), 31--38.
PMID: 24149292; PMCID: PMC3737890.
\url{https://pmc.ncbi.nlm.nih.gov/articles/PMC3737890/}

% [29]
\bibitem{Lee2005}
Lee, C.-H., Lee, Y.-J. \& Cheong, C.-C.
\textit{A Kinematical Analysis of the Taekwondo Apchagi}.
In: Proceedings of the 23rd International Symposium on Biomechanics in Sports (ISBS), Beijing, China, 595--597, 2005.
\bibitem{Echeverria2021}
Echeverria, J. \& Santos, O.~C.
\textit{Toward Modeling Psychomotor Performance in Karate Combats
Using Computer Vision Pose Estimation}.
Sensors, 21(24): 8378, 2021.
\href{https://doi.org/10.3390/s21248378}{https://doi.org/10.3390/s21248378}
\bibitem{OpenPoseLicense}
Carnegie Mellon University,
\textit{OpenPose: Academic or Non-Profit Organization Noncommercial
Research Use Only License},
Online verfügbar unter:
\url{https://github.com/CMU-Perceptual-Computing-Lab/openpose/blob/master/LICENSE},
abgerufen am 6. September 2026.
\bibitem{MediaPipeLicense}
Google LLC,
\textit{MediaPipe, Apache License Version 2.0},
Online verfügbar unter:
\url{https://github.com/google-ai-edge/mediapipe/blob/master/LICENSE},
abgerufen am 6. September 2026.

\bibitem{Brooke1996}
Brooke, J.
\textit{SUS: A Quick and Dirty Usability Scale}.
In: Jordan, P. W., Thomas, B., McClelland, I. L. \& Weerdmeester, B. (Hrsg.),
\textit{Usability Evaluation in Industry}, Taylor \& Francis, 189--194, 1996.

\bibitem{SauroLewis2016}
Sauro, J. \& Lewis, J. R.
\textit{Quantifying the User Experience: Practical Statistics for User Research}.
2. Auflage, Morgan Kaufmann, 2016.

\end{thebibliography}

\end{document}
